\pdfoutput=1
\documentclass[11pt]{amsart}
\usepackage{amsmath,amssymb,amsthm}
\usepackage{mathtools}
\usepackage[margin=1.2in]{geometry}
\usepackage{booktabs}
\usepackage{array}
\usepackage{flafter}
\usepackage{xcolor}
\usepackage[colorlinks=true,linkcolor=blue,citecolor=blue,urlcolor=blue]{hyperref}
\hypersetup{pdfauthor={Bernd Johannes Wuebben},
  pdftitle={A vanishing-cycle obstruction to smoothing Calabi-Yau threefolds}}

\setlength{\emergencystretch}{2em}

\newtheorem{theorem}{Theorem}[section]
\newtheorem{proposition}[theorem]{Proposition}
\newtheorem{lemma}[theorem]{Lemma}
\newtheorem{corollary}[theorem]{Corollary}
\newtheorem{conjecture}[theorem]{Conjecture}
\theoremstyle{definition}
\newtheorem{definition}[theorem]{Definition}
\newtheorem{example}[theorem]{Example}
\newtheorem{question}[theorem]{Question}
\newtheorem{convention}[theorem]{Convention}
\theoremstyle{remark}
\newtheorem{remark}[theorem]{Remark}

\newcommand{\ZZ}{\mathbb{Z}}
\newcommand{\QQ}{\mathbb{Q}}
\newcommand{\RR}{\mathbb{R}}
\newcommand{\CC}{\mathbb{C}}
\newcommand{\PP}{\mathbb{P}}
\newcommand{\cO}{\mathcal{O}}
\newcommand{\cX}{\mathcal{X}}
\newcommand{\conv}{\operatorname{conv}}
\newcommand{\Cone}{\operatorname{Cone}}
\newcommand{\Sing}{\operatorname{Sing}}
\newcommand{\Pic}{\operatorname{Pic}}
\newcommand{\Cl}{\operatorname{Cl}}
\newcommand{\Def}{\operatorname{Def}}
\newcommand{\rk}{\operatorname{rank}}
\newcommand{\Bl}{\operatorname{Bl}}
\newcommand{\dP}{\mathrm{dP}}
\newcommand{\GL}{\mathrm{GL}}
\newcommand{\Xh}{\widehat{X}}
\newcommand{\Sb}{\bar{S}}
\newcommand{\Disk}{\mathbb{D}}

\title[A vanishing-cycle smoothing obstruction]{A vanishing-cycle
obstruction to smoothing\\ Calabi--Yau threefolds}
\author{Bernd Johannes Wuebben}
\subjclass[2020]{14J32 (primary); 14B07, 14M25, 32S30, 32S55 (secondary)}
\keywords{Calabi--Yau threefolds, smoothing of singularities, vanishing
cycles, del Pezzo cones, ordinary double points, Milnor fibers,
reflexive polytopes}

\begin{document}

\begin{abstract}
A projective Calabi--Yau threefold whose singular points all admit
local smoothings need not admit a global one.  For threefolds whose
singularities are ordinary double points and anticanonical cones over
del Pezzo surfaces of degrees~$6$ and~$7$, we prove a necessity
criterion: any one-parameter smoothing produces a rational homology
class on the disjoint union of the singularity links that vanishes on
a crepant resolution, lies in a canonical root subspace of
$K^{\perp}$ at each del Pezzo cone point, and is nonzero at every
node, at every degree-$7$ point, and at every degree-$6$ point
smoothed through its one-parameter component.  The engine is a period
computation: the integral of the holomorphic $3$-form over a
vanishing cycle equals the local smoothing parameter times a unit, so
a smoothed germ cannot have rationally trivial vanishing homology.
In the purely nodal case the criterion recovers the necessity half of
Friedman's relation for arbitrary one-parameter smoothings, without
appeal to unobstructedness.  The criterion is a finite linear-algebra
test, and it obstructs: the $\Delta$-regular anticanonical
hypersurface of the unique doubly isolated Batyrev mirror pair, with
$26$ nodes, two degree-$7$ and two degree-$6$ cone points, and every
germ locally smoothable, admits no smoothing.  The same argument
applies to deformations that leave a germ with no smoothing in place,
and shows that the mirror partner of that hypersurface, already
non-smoothable because one of its germs admits no smoothing, carries
the global obstruction as well.  It also obstructs a
hypersurface with $17$ nodes and one degree-$7$ point whose nodal
configuration satisfies the classical relation condition, so that the
nodal theory alone finds nothing wrong.  On a further hypersurface, with $14$ nodes and one
degree-$7$ point, the test is silent; we pose its smoothability, a
genuinely mixed nodal and divisorial smoothing if it exists, as an
open problem.
\end{abstract}

\maketitle

\section{Introduction}\label{sec:intro}

\subsection{The problem}
Let $X$ be a projective threefold over $\CC$ with trivial dualizing
sheaf and isolated singularities.  When every singular germ of $X$
admits a smoothing, one may ask whether $X$ itself does.  For ordinary
double points the answer is classical: by Friedman's
criterion~\cite{Friedman86}, simultaneous smoothing is governed by a
relation among the exceptional curves of a small resolution in which
every curve appears with nonzero coefficient; see also
Rollenske--Thomas~\cite{RT09} for a topological formulation in all odd
dimensions.  Outside the nodal case much less is known.  The local
deformation theory of the simplest non-nodal germs, the
anticanonical cones over del Pezzo surfaces, has been understood
since Altmann~\cite{Altmann95,Altmann97}, and their smoothing
components have explicit geometric descriptions
(Gross~\cite[\S5.5]{Gross97}; Petracci~\cite{Petracci20,Petracci19}).
What has been missing is a global criterion that couples nodal and
divisorial smoothing data.

The first three papers of this
program~\cite{PaperI,PaperII,PaperIII} classified the isolated
Gorenstein toric threefold germs arising on Batyrev hypersurfaces from
unit-edge polygons, constructed global smoothings when a
dual-edge-length invariant is at least two, and isolated a unique
Batyrev mirror pair, associated with a $22$-vertex reflexive
$4$-polytope $\Delta_{\mathbb F_1}$ and its polar, for which both
members have only isolated singularities and the remaining smoothing
question is genuinely mixed: one member carries a cone singularity
with no smoothing and is locally obstructed, while the other, denoted
$X^{\circ}$ throughout, has $26$ ordinary double points, two
anticanonical cones over the del Pezzo surface of degree~$7$, and two
over the del Pezzo surface of degree~$6$: thirty germs, every one of
them locally smoothable.  A purely nodal analysis cannot decide
$X^{\circ}$: the node relation matrix has two coloops, so any smoothing
would have to recruit the divisorial deformation directions.

\subsection{Results}
Fix the following setting.  $X$ is a compact complex threefold with
$\omega_X\cong\cO_X$ and $h^1(\cO_X)=0$ whose singular points are
ordinary double points and anticanonical cones over del Pezzo surfaces
of degrees~$6$ and~$7$; $\pi\colon\Xh\to X$ is a crepant resolution,
small over the nodes and with an exceptional del Pezzo surface $E_p$
over each cone point $p$.  For a singular point $p$ let $L_p$ denote
its link.  A \emph{one-parameter deformation} of $X$ is a flat proper
analytic family $\cX\to\Disk$ over a disk $\Disk$ with
$\cX_0\cong X$, and it is a \emph{one-parameter smoothing} if $X_s$ is
smooth for all $s\neq0$.  Such a smoothing induces at each singular
point a smoothing of the local germ, with Milnor fiber $M_p$, and we
set
\[
R_p\ :=\ \ker\bigl(H_2(L_p;\QQ)\longrightarrow H_2(M_p;\QQ)\bigr),
\]
the part of the link homology that bounds in the local Milnor fiber.

\begin{theorem}[mixed necessity; Theorem~\ref{thm:necessity}]
\label{thm:B-intro}
Suppose $X$ as above admits a one-parameter smoothing.  Then the space
\[
K(X)\ :=\ \Bigl\{\kappa\in\ker\Bigl(\bigoplus_pH_2(L_p;\QQ)
\to H_2(\Xh;\QQ)\Bigr)\ :\ \kappa_p\in R_p
\text{ at every cone point }p\Bigr\}
\]
contains, for every node $q$, for every degree-$7$ cone point $q$, and
for every degree-$6$ cone point $q$ whose induced local smoothing lies
in the one-parameter component of its miniversal base, a class
$\kappa$ with $\kappa_q\neq0$.  Here the maps are induced by the
inclusions $L_p\subset X\smallsetminus\Sing X\subset\Xh$.
\end{theorem}

The subspaces $R_p$ are computed in closed form
(Proposition~\ref{prop:kernels}): at a degree-$7$ point $R_p$ is the
unique line spanned by a $(-2)$-class of
$K^{\perp}\subset H_2(E_p;\QQ)$, and at a degree-$6$ point it is the
$A_1$ or the $A_2$ summand of the root lattice
$K^{\perp}\cong A_1\oplus A_2$, according to the component of the local
miniversal base through which the family smooths the germ.  These
subspaces coincide with the lattices $\Xi(V)$ attached by
Kuznetsov--Prokhorov~\cite[Theorem~1.9]{KP23} to the del Pezzo
threefolds $V$ that sweep out the local smoothings.

Theorem~\ref{thm:B-intro} is a finite linear-algebra test on any
concrete $X$, and on the distinguished hypersurface it obstructs.

\begin{theorem}[Theorem~\ref{thm:main}]\label{thm:A-intro}
Let $X^{\circ}$ be a $\Delta$-regular anticanonical hypersurface in
the Gorenstein Fano toric fourfold associated with the polar of
$\Delta_{\mathbb F_1}$.  Then $X^{\circ}$ admits no one-parameter
smoothing.  Moreover no fiber of a small representative of the
miniversal deformation of $X^{\circ}$ is smooth, and no formal
arc in the miniversal deformation of $X^{\circ}$ meets all thirty
local smoothing loci.
\end{theorem}

The mirror partner $X_{\Delta}$ of $X^{\circ}$ in the doubly isolated
pair carries an anticanonical cone over $\mathbb F_1$, a germ with no
smoothing~\cite{PaperI,PaperIII}, so both members of the pair are
non-smoothable.  The criterion says more.  Extended to deformations
that leave such a germ in place (Section~\ref{sec:partner}), it shows
that $X_{\Delta}$ carries the global obstruction as well: two of its
twenty nodes are smoothed by no one-parameter deformation at all, and
this is
independent of the $\mathbb F_1$ germ, which already prevents a
smoothing for
a purely local reason.  The two members of the pair are therefore not
an instance of one obstruction seen locally and globally; each fails
globally, and one of them fails locally in addition.

Two by-products deserve separate statement.  First, the mechanism
behind Theorem~\ref{thm:B-intro} is elementary and quantitative: the
period of the holomorphic $3$-form of the nearby fiber over the
vanishing cycle of a node equals the local smoothing parameter times a
unit (Lemma~\ref{lem:P}), and an analogous nonvanishing holds for the
distinguished cycles of the degree-$7$ and one-parameter degree-$6$
smoothings (Lemma~\ref{lem:D}), where the decisive constant is an
explicit period equal to $4\pi^2$ on a cycle we construct by hand.
Consequently a smoothed germ cannot have rationally null-homologous
vanishing homology, at any order of tangency of the family with the
local discriminants.  In the purely nodal case this yields:

\begin{corollary}[Corollary~\ref{cor:friedman}]\label{cor:friedman-intro}
For a compact complex threefold with $\omega\cong\cO$,
$h^1(\cO)=0$ and only ordinary double points, any one-parameter
smoothing forces, for every node $q$, the existence of a relation
among the exceptional curve classes of a small resolution that is
nonzero at $q$.  This is the necessity half of Friedman's criterion,
for arbitrary one-parameter smoothings, obtained without
unobstructedness of the deformation space.
\end{corollary}

Second, on $X^{\circ}$ the argument localizes the obstruction
topologically:

\begin{proposition}[Proposition~\ref{prop:coloops}]
\label{prop:coloops-intro}
In any hypothetical one-parameter smoothing of $X^{\circ}$, the
vanishing $3$-spheres of the two coloop nodes would be null-homologous
in $H_3(X_s;\QQ)$.
\end{proposition}

The two statements collide (Lemma~\ref{lem:P} forbids exactly this),
and that collision is the proof of Theorem~\ref{thm:A-intro}.  A
complementary second proof, with a different period input, runs
through a degree-$7$ point instead of
the nodes (Section~\ref{sec:proofs}).

Finally, the criterion is sharp in the sense that it does not obstruct
indiscriminately.  Applying the same finite test to the anticanonical
hypersurface $X_{19}$ of a $19$-vertex reflexive polytope from the
companion paper~\cite{PaperIII} (fourteen nodes and one degree-$7$
cone point, again with two nodal coloops), the finite test is
silent: no covered coordinate is forced to vanish
(Section~\ref{sec:test}), and we ask:

\begin{question}\label{q:X19}
Is $X_{19}$ smoothable?  Any smoothing of $X_{19}$ would be a
simultaneous smoothing of nodal and del Pezzo cone singularities that
no purely nodal mechanism can produce.
\end{question}

\subsection{Method}
The proof has three independent inputs.  \emph{Membership}: a surgery
decomposition of the nearby fiber along the links, together with
Mayer--Vietoris, places the boundary of every $3$-cycle of $X_s$ into
the space $K(X)$ of Theorem~\ref{thm:B-intro}; the computation of the
subspaces $R_p$ rests on the classification of del Pezzo threefolds of
degrees~$6$ and~$7$ (Fujita~\cite{Fujita80a}) and on an identification,
apparently not available in the literature in the form needed, of the
germ Milnor fiber with a del Pezzo threefold complement \emph{as a
manifold with boundary}, compatibly with the two natural circle
fibrations of the link (Lemma~\ref{lem:boundary}).  \emph{Periods}:
the vanishing-period computations of Section~\ref{sec:periods} convert
nonvanishing of local smoothing parameters into nonvanishing of
homology classes.  \emph{Finite input}: for $X^{\circ}$, an exact
integral computation on a fixed crepant resolution shows that on the
relevant restricted kernels the coordinates of the two coloop nodes
and of the degree-$7$ directions vanish identically
(Section~\ref{sec:example}); all decisive matrix data are printed, and
every computation is reproduced by the ancillary programs
(Appendix~\ref{app:cert}).

We emphasize what the argument does \emph{not} use.  It does not use
the first-order deformation theory of $X$ beyond motivation: the
localization map, its image, and the comparison of coordinate systems
on local deformation spaces never enter.  It is insensitive to the
order of tangency of the smoothing family with the local
discriminants, the phenomenon that blocks first-order methods: the
necessity statements of \cite{Friedman86,RT09} concern the first-order
data of a deformation class, and their proofs are specific to nodes
(one-dimensional local $T^1$, quadric exceptional geometry), whereas
the periods above see the actual arc.  Neither does it use
unobstructedness, which is unavailable for the non-lci cone germs.

\subsection{Relation to prior work}
Friedman~\cite{Friedman86} proved the nodal criterion;
Rollenske--Thomas~\cite{RT09} gave a topological proof and the
higher-dimensional statement, in both cases for first-order data, with
the passage to actual smoothings supplied by unobstructedness
(Kawamata, Tian, Ran) in the nodal case.  Friedman--Laza
\cite{FL22a,FL22b} identify the first-order localization data of
threefolds with isolated singularities in terms of link homology; our
Section~\ref{sec:example} consumes only the resolution-side avatar of
that picture, in the form of an integral matrix.
Gross~\cite{Gross97} determined the smoothing components of the del
Pezzo cone germs; Altmann~\cite{Altmann95,Altmann97} computed their
miniversal deformations; Petracci~\cite{Petracci20,Petracci19} and
Kaloghiros--Petracci~\cite{KP21} give modern treatments, the latter
computing the Betti numbers of the degree-$6$ Milnor fibers, in
agreement with Section~\ref{sec:germs}.  The classification of the
sweeping del Pezzo threefolds is due to Fujita~\cite{Fujita80a}; we
follow the caution recorded there (Supplementary notes) that the
degree-$6$ case of Iskovskikh's earlier list omits
$\PP(T_{\PP^2})$, and cite Fujita, with
Kuznetsov--Prokhorov~\cite{KP23} as the modern reference.  On the
constructive side, multi-parameter deformations of toric varieties in
a single character degree were built by
Ilten--Vollmert~\cite{IV12}; the sufficiency counterpart of
Theorem~\ref{thm:B-intro}, in particular Question~\ref{q:X19}, is the
subject of work in progress and is not addressed here.

\begin{convention}\label{conv:dp}
$\dP_d$ denotes a del Pezzo \emph{surface} of degree $d=K^2$; thus
$\dP_6=\Bl_3\PP^2$ and $\dP_7=\Bl_2\PP^2$.  A \emph{del Pezzo
threefold} is a smooth projective threefold $V$ with $-K_V=2H$ for an
ample $H$; its degree is $H^3$.  For a polygon or polytope, edges are
\emph{unit} if they have lattice length one.  Homology and cohomology
have $\QQ$-coefficients unless stated otherwise.  The anticanonical
cone over a del Pezzo surface $S$ is
$C(S)=\operatorname{Spec}\bigoplus_{k\ge0}H^0(S,-kK_S)$.
\end{convention}

\section{The germs: links, miniversal bases, Milnor fibers}
\label{sec:germs}

Throughout this section $(Y,0)=C(S)$ is the anticanonical cone over
a smooth del Pezzo surface $S$, an isolated Gorenstein toric threefold
singularity when $S$ is toric.  For $S=\dP_6$, $\dP_7$ and
$\mathbb F_1$, the cases used in this paper, it is the toric variety
of the cone over the unit-edge reflexive hexagon, pentagon and
quadrilateral respectively, placed at height one.  (The
unit-edge polygon here is the \emph{fan} polygon of $S$; its polar,
the Newton polygon of $-K_S$, is not unit-edge for the pentagon.)

\subsection{Links}\label{sec:links}
The punctured cone is the total space of $\cO_S(K_S)$ minus its zero
section, so the link $L$ is the unit circle bundle of $\cO_S(K_S)$,
and the homology Gysin sequence together with $H_1(S)=0$ gives a
canonical isomorphism
\begin{equation}\label{eq:linkH2}
H_2(L;\QQ)\ \cong\ K_S^{\perp}\ \subset\ H_2(S;\QQ),
\end{equation}
of rank $3$ for $\dP_6$ and $2$ for $\dP_7$.  For a node,
$L\cong S^2\times S^3$ and $H_2(L;\QQ)=\QQ$.  The lattice
$K^{\perp}(\dP_6)$ is the root lattice $A_1\oplus A_2$ (discriminant
$6$); $K^{\perp}(\dP_7)$ has discriminant $7$ and contains exactly one
pair $\pm$ of primitive classes of square $-2$.  Both facts are
elementary and are certified by the ancillary programs together with
everything else in this section that is finite.

\subsection{Miniversal deformations}\label{sec:miniversal}
By Altmann's theorems, the deformation theory of these germs is
completely explicit.  $T^1$ is concentrated in the Gorenstein
character degree, of dimension $N-3$ for an $N$-gon
\cite[(6.5) and the subsequent Remark]{Altmann95}; the versal family
in that degree, constructed in \cite{Altmann97}, is therefore a versal
deformation of the affine cone, with base ring computed there:
\[
\begin{aligned}
\dP_7&:\ \ \CC[s_1,s_2]/(s_1^2,\,s_1s_2)
&&\text{\cite[(9.1)(v)]{Altmann97}},\\
\dP_6&:\ \ \CC[s_1,s_2,s_3]/(s_1s_3,\,s_2s_3)
&&\text{\cite[(2.5)]{Altmann97}}.
\end{aligned}
\]
The reduced components correspond to the maximal decompositions of the
polygon into Minkowski sums of lattice polygons
\cite[(2.5),\,(8.1)]{Altmann97}: the pentagon decomposes in exactly
one way, into a unit segment pair and a unimodular triangle, giving
one reduced component, the line $\{s_1=0\}$; the hexagon decomposes as
a sum of two triangles (the line $\{s_1=s_2=0\}$) and as a sum of
three segments (the plane $\{s_3=0\}$).

Three standard supplements, used later, deserve explicit statement.
First, \emph{minimality}: the Kodaira--Spencer map of Altmann's family
is an isomorphism onto $T^1$ \cite[(6.2),\,(6.3)]{Altmann97}, so the
versal family is miniversal.  Second, \emph{the germ}: $T^1$ and $T^2$
of the affine cone are finite-dimensional and supported at the vertex,
and the deformation theory of the affine variety is identified with
that of the analytic germ $(Y,0)$ by Artin
approximation~\cite{Artin69}; the analytic miniversal deformation
exists by Grauert~\cite{Grauert72}; the base is the analytic germ of
the ring displayed above.  (The versality criterion used in
\cite[(7.2)]{Altmann97} is that of de~Jong--van~Straten
\cite[\S4]{dJvS94}.)  Third, \emph{arcs}: a morphism from the germ of
a disk to the miniversal base factors through the reduction and,
since the analytic local ring of a disk is a domain, lies in a single
reduced component; for the pentagon base this forces $s_1\equiv0$
along any arc, and for the hexagon base it forces the arc into the
line or into the plane.  We refer to the corresponding local
smoothings as smoothings \emph{through} the given component.

\subsection{Sweeps and Milnor fibers}\label{sec:sweeps}
The smoothing components of these cones are cones swept by hyperplane
sections, in the sense of Pinkham
\cite[(4.2),\,(5.1),\,(5.5)]{Pinkham74}; see
\cite[Lemma~4.1, Proposition~4.2]{PP15} for the statement in all
dimensions in the pencil form used here.  Concretely
(Gross~\cite[\S5.5]{Gross97};
Petracci~\cite[Proposition~2.2]{Petracci20},
\cite[Proposition~4.12]{Petracci19}): the general fiber of the
degree-$7$ smoothing is $V_7\smallsetminus\Sb$ with
$V_7=\Bl_{pt}\PP^3$; the general fiber of the hexagon's one-parameter
(line) component is $(\PP^1)^3\smallsetminus\Sb$; and that of its
two-parameter (plane) component is
$W\smallsetminus\Sb$ with $W=\PP(T_{\PP^2})$ the divisor of bidegree
$(1,1)$ in $\PP^2\times\PP^2$.  In each case $V$ is a del Pezzo
threefold of the same degree with $\Sb\in|H|$ a copy of $S$, embedded
by $-K_S=H|_{\Sb}$; the classification of the possible $V$
(Fujita~\cite[(5.8),\,(5.16)]{Fujita80a}) shows there are no others.
For the record, the open Milnor fibers of the three channels have
$(b_2,b_3)=(1,1)$, $(2,1)$ and $(1,2)$ respectively; the degree-$6$
values agree with \cite[Proposition~3.5]{KP21}.

\begin{proposition}\label{prop:kernels}
Let $M=V\smallsetminus\Sb$ be one of the three sweep complements.
Under the identification~\eqref{eq:linkH2} of $H_2$ of the link
$L=\partial(\text{tube around }\Sb)$ with $K^{\perp}_{S}$,
\[
\ker\bigl(H_2(L;\QQ)\to H_2(M;\QQ)\bigr)
=\ker\bigl(K^{\perp}_S\to H_2(V;\QQ)\bigr)=
\begin{cases}
\text{the }(-2)\text{-line} & V=V_7,\\
A_1 & V=(\PP^1)^3,\\
A_2 & V=W,
\end{cases}
\]
where the map is the restriction to $K^\perp$ of the pushforward
$H_2(S;\QQ)\to H_2(V;\QQ)$.  For a node,
$\ker(H_2(L;\QQ)\to H_2(M;\QQ))=H_2(L;\QQ)=\QQ$.
\end{proposition}

\begin{proof}
Since $\Sb$ is an ample divisor with $H_1(\Sb)=0$, the Thom
isomorphism gives $H_3(V,M)\cong H_1(\Sb)=0$, so
$H_2(M)\hookrightarrow H_2(V)$, and the inclusion $L\subset M$
composed with this injection factors, up to homotopy through the
tube, as the projection $L\to\Sb$ followed by $\Sb\subset V$; the two
kernels therefore coincide.  The pushforward is computed in the
standard bases.  For $V_7$: a smooth member of $|H|=|2H_{\PP^3}-E|$
is the strict transform of a quadric through the blown-up point, so
$\Sb\cong\Bl_2\PP^2$ with $H_{\PP^3}|_{\Sb}=2\ell-e_1-e_2$ and
$E|_{\Sb}=\ell-e_1-e_2$; the pushforward matrix has rows
$(2,1),(1,1),(1,1)$ on $(\ell,e_1,e_2)$ and kernel
$\langle e_1-e_2\rangle$, which is the $(-2)$-line.  For
$(\PP^1)^3$: the three projections restrict to the conic pencils
$|\ell-e_i|$, the rows are $(1,1,1)$ and the unit vectors, and the
kernel is $\langle\ell-e_1-e_2-e_3\rangle=A_1$.  For $W$: the two
contractions restrict to $\ell$ and $2\ell-e_1-e_2-e_3$, and the
kernel is $\langle e_1-e_2,\,e_2-e_3\rangle=A_2$.  For the node,
$M\simeq T^*S^3$ has $H_2=0$.
\end{proof}

\begin{remark}\label{rem:xi}
The three kernels are the lattices $\Xi(V)$ of
Kuznetsov--Prokhorov: \cite[Theorem~1.9(i)]{KP23} states
$\Cl(V)\cong\Xi(V)^{\perp}\subset\Cl(\Sb)$ with $\Xi(V)$ negative
definite of rank $10-d-\rho(V)$, giving $A_1$, $A_1$, $A_2$ for
$V_7$, $(\PP^1)^3$, $W$, an independent confirmation.
\end{remark}

\subsection{Which sweep for which component}\label{sec:match}
The one-parameter (line) component of the hexagon base sweeps to
$(\PP^1)^3$ and the two-parameter (plane) component to $W$: this is
stated in \cite[Proposition~2.2]{Petracci20}, follows from
\cite[(7.1.5)]{Altmann95} (the total spaces over the two components are
the cones over $(\PP^1)^3$ and over $\PP^2\times\PP^2$, whose
anticanonical hyperplane slice is $W$), and is confirmed numerically
by the Betti numbers quoted above against
\cite[Proposition~3.5]{KP21}.  The results of
Sections~\ref{sec:periods}--\ref{sec:proofs} are stated so that only
the degree-$7$ channel and the fact that each degree-$6$ kernel is
$A_1$ or $A_2$ are consumed by the main theorem; the matching enters
only in the sharper form of Theorem~\ref{thm:necessity}(3).

\section{Vanishing periods}\label{sec:periods}

This section contains the analytic mechanism.  Throughout,
$\cX\to\Disk$ is a one-parameter smoothing of a compact complex
threefold $X$ with $\omega_X\cong\cO_X$, $h^1(\cO_X)=0$ and isolated
singularities of the types in Theorem~\ref{thm:B-intro}.

\subsection{A relative holomorphic \texorpdfstring{$3$}{3}-form}\label{sec:relform}
Flatness and the Gorenstein property of all fibers make
$\cX\to\Disk$ a Gorenstein morphism after shrinking the disk, so the
relative dualizing sheaf $\omega_{\cX/\Disk}$ is invertible and
compatible with base change \cite{Conrad00}.  We first prove
$c_1(\omega_{X_s})=0$ integrally.  Choose the disjoint good
representatives and exterior $W$ of Section~\ref{sec:decomp}.  Over a
neighborhood of $W$ the family is a smooth proper submersion with
boundary and hence is $C^\infty$-trivial.  Since
$\omega_X|_W\cong\cO_W$, the restriction of
$c_1(\omega_{\cX/\Disk})$ to the corresponding exterior in $X_s$
vanishes.  On each local piece it also vanishes: choose first a Stein
contractible neighborhood $N_p$ of $(p,0)$ in $\cX$, on which the
invertible sheaf $\omega_{\cX/\Disk}$ is trivial, and then the good
representative $U_p$ inside $N_p$.

For every boundary link in this paper, $H^1(L_p;\ZZ)=0$.  This is
immediate for the nodal link $S^2\times S^3$.  For a cone point, $L_p$
is the circle bundle of $\cO_S(K_S)$ over the simply connected del
Pezzo surface $S$; the homotopy sequence shows that $\pi_1(L_p)$ is
the cyclic group $\ZZ/d$, where $d$ is the divisibility of the Euler
class $K_S\neq0$, because the connecting map
$\pi_2(S)\to\pi_1(S^1)=\ZZ$ is evaluation against $K_S$.  Hence
$H^1(L_p;\ZZ)=\operatorname{Hom}(\ZZ/d,\ZZ)=0$ for every del Pezzo
cone point, whether or not $K_S$ is primitive.  (It is primitive for
$\dP_6$, $\dP_7$ and $\mathbb F_1$, where $L_p$ is simply connected,
but not for $\PP^1\times\PP^1$.)  The Mayer--Vietoris
sequence for $X_s=W\cup\bigsqcup_pM_p$ therefore injects
\[
 H^2(X_s;\ZZ)\lhook\joinrel\longrightarrow
 H^2(W;\ZZ)\oplus\bigoplus_pH^2(M_p;\ZZ).
\]
The preceding local vanishings give $c_1(\omega_{X_s})=0$.
Semicontinuity gives $h^1(\cO_{X_s})=0$; the exponential sequence then
embeds $\Pic(X_s)$ into $H^2(X_s;\ZZ)$ and forces
$\omega_{X_s}\cong\cO_{X_s}$.  (The hypothesis
$\omega_X\cong\cO_X$ cannot be dropped: the cone over a smooth quadric
surface in $\PP^4$ is a one-node threefold smoothed by the family of
quadric threefolds, and the vanishing cycle of that smoothing
\emph{is} null-homologous; that example satisfies $h^1(\cO_X)=0$ and
is excluded here exactly because $\omega_X\cong\cO_X(-3)$ is
nontrivial.  The hypothesis $h^1(\cO_X)=0$ enters only through the
exponential sequence in the previous sentence.)  Constants and upper
semicontinuity give $h^0(\cO_{X_s})=1$ after shrinking, hence
$h^0(\omega_{X_s})=1$.  By Grauert's base-change theorem
\cite{Grauert60}, $\pi_*\omega_{\cX/\Disk}$ is then a line bundle on the shrunk
disk; we fix a trivializing section $\Omega$, whose restriction
$\Omega_s$ generates the invertible sheaf $\omega_{X_s}$ and is
therefore a nowhere-vanishing holomorphic $3$-form on the smooth locus
of $X_s$ for $s\neq0$ and a generator of $\omega_X$ near each singular point at
$s=0$.

\subsection{Nodes}\label{sec:nodeperiod}
\begin{lemma}\label{lem:P}
Let $p$ be an ordinary double point of $X$.  Then the vanishing-cycle
class $\delta_p(s)\in H_3(X_s;\QQ)$ is nonzero for all small
$s\neq0$.
\end{lemma}

\begin{proof}
The germ $(X,p)$ is $\{xy-zw=0\}\subset(\CC^4,0)$, with miniversal
deformation $\{xy-zw=u\}$ over the $u$-line
\cite{KasSchlessinger72}.  The family induces a classifying map
$s\mapsto\tilde u(s)$, holomorphic with $\tilde u(0)=0$ (only its
nonvanishing at each fixed $s\neq0$ is used, which is independent of
choices), and $X_s$
is smooth near $p$ exactly when $\tilde u(s)\neq0$; the hypothesis
gives $\tilde u(s)\neq0$ for all small $s\neq0$.  The pullback total
space is the hypersurface $\{xy-zw=\tilde u(s)\}$, Gorenstein even at
its singular points, and adjunction exhibits $\omega_{\cX/\Disk}$
near $p$ as generated by the relative Gelfand--Leray form
$\GL_u=\operatorname{Res}\bigl[dx\,dy\,dz\,dw/(xy-zw-u)\bigr]$.
Hence $\Omega=g\cdot\GL$ with $g$ holomorphic on the total-space germ
and $g(p,0)\neq0$; extend $g$ to an ambient holomorphic function so
that $g(0,s)$ is defined.

After the linear change putting $xy-zw$ into the form
$\sum_{i=1}^4v_i^2$, the vanishing cycle over $u$ is the real slice
$\delta_u=\sqrt u\cdot S^3$, whose class is single-valued: in fiber
dimension three the Picard--Lefschetz transvection fixes $\delta$
(as $\langle\delta,\delta\rangle=0$), and the antipodal ambiguity of
$\sqrt u$ acts on $S^3$ with degree $+1$.  Under
$v\mapsto\lambda v$ one has $u\mapsto\lambda^2u$ and
$\GL\mapsto\lambda^2\GL$, so the local period is exactly linear:
\[
\int_{\delta_u}\GL_u\;=\;c_0\,u,
\qquad
c_0=\int_{S^3}\frac{dv_1\,dv_2\,dv_3}{2v_4}=\pi^2\neq0
\]
up to orientation conventions.  For the twisted form, expanding the
ambient extension $g=g(p,s)+O(|v|)$ and using $|v|=|u|^{1/2}$ on
$\delta_u$ gives
\[
F(u,s):=\int_{\delta_u}g\cdot\GL_u
= u\,\bigl(c_0\,g(p,s)+O(|u|^{1/2})\bigr),
\]
so $F(u,s)=u\cdot(\text{unit})$ near $(0,0)$; here $u,s$ are
independent until the substitution $u=\tilde u(s)$.

Suppose $\delta_p(s_0)=0$ in $H_3(X_{s_0};\QQ)$ for some small
$s_0\neq0$.  A nonzero integer multiple of the cycle then bounds, and
$\Omega_{s_0}$ is closed, so
$\int_{\delta_p(s_0)}\Omega_{s_0}=0$.  But this period equals
$F(\tilde u(s_0),s_0)=\tilde u(s_0)\cdot(\text{unit})\neq0$, a
contradiction.
\end{proof}

\subsection{Del Pezzo cone points}\label{sec:dpperiod}
For the cone germs the same strategy applies, with the local model
family replaced by the sweep and the explicit period computed on the
sweep complement.  Both relevant complements admit a uniform toric
description.  Let $V$ be $(\PP^1)^3$ or $\Bl_{pt}\PP^3$ with its fan
in $N\cong\ZZ^3$, and let $m\in M$ be the character $(1,1,1)$,
respectively $(1,1,-1)$ in the basis where the fan of
$\Bl_{pt}\PP^3$ has rays $e_1,e_2,e_3$, $r=-e_1-e_2-e_3$ and
$e_1+e_2+e_3$.  In both cases $\langle m,v\rangle=\pm1$ on every ray,
the closure $\Sb$ of $\{\chi^m=1\}$ is a copy of the appropriate del
Pezzo surface with $[\Sb]=-\tfrac12K_V$, and
\begin{equation}\label{eq:omegaM}
\Omega_M=\frac{w}{(w-1)^2}\,
\frac{dx}{x}\wedge\frac{dy}{y}\wedge\frac{dz}{z},
\qquad w=\chi^m,
\end{equation}
is the generator of $\omega(V\smallsetminus\Sb)\cong\cO$, unique up to
scale: the order of \eqref{eq:omegaM} along a boundary divisor $D_v$
is $\langle m,v\rangle-2\min(\langle m,v\rangle,0)-1=0$, and along
$\Sb$ it is $-2$.  Here uniqueness is algebraic.  If
$f\in\cO(V\smallsetminus\Sb)^\times$, then
$\operatorname{div}_V(f)=n\Sb$ for some $n\in\ZZ$; since
$\Pic(V)$ is torsion-free and $[\Sb]\neq0$, one has $n=0$, and a
rational function without zeros or poles on the projective variety
$V$ is constant.  (These statements, the slice combinatorics below,
and the constancy claim in the next proof are certified by the
ancillary programs.)

\begin{lemma}\label{lem:D1}
In each of the two complements $M=V\smallsetminus\Sb$ above there is a
closed $3$-cycle $T\subset M$ with $\int_T\Omega_M=\pm4\pi^2$.
Consequently $[\Omega_M]\neq0$ in $H^3(M;\CC)$, and since
$b_3(M)=1$, the period pairing
$H_3(M;\QQ)\to\CC$, $\nu\mapsto\int_\nu\Omega_M$, is injective.
\end{lemma}

\begin{proof}
Choose an edge $\langle v_+,v_-\rangle$ of the fan with
$\langle m,v_\pm\rangle=\pm1$ and an adjacent maximal cone
$\langle v_+,v_-,n_0\rangle$ with $\gamma:=\langle m,n_0\rangle=+1$;
such pairs exist in both cases, any admissible choice works, and the
transitions printed below arise from
$(v_+,v_-,n_0;n_1)=(e_1,-e_2,e_3;-e_3)$ for $(\PP^1)^3$ and
$(e_1,r,e_2;e_3)$ for $\Bl_{pt}\PP^3$.  In the unimodular chart with coordinates
$(p,q,z)$ dual to $(v_+,v_-,n_0)$ one has $w=pq^{-1}z^{\gamma}$,
$\Sb=\{t=0\}$ with $t=q-pz^{\gamma}$, and
$\Omega_M=\pm z^{\gamma-1}\,dp\wedge dq\wedge dz/t^{2}$.
Fix $\varepsilon>0$, $R>1$ and set
\begin{align*}
T_1&=\{p=0,\ q=\varepsilon R^{\gamma}e^{i\psi},\ |z|\le R\},\\
T_2&=\{p=r\varepsilon e^{i\theta},\
q=-(1-r)\varepsilon R^{\gamma}e^{i(\theta+\gamma\varphi)},\
z=Re^{i\varphi}\},\qquad r\in[0,1],\\
T_3&=\{p'=\varepsilon R^{c}e^{i\psi},\ q'=0,\ |z'|\le 1/R\}
\quad\text{in the chart at the opposite corner }n_1,
\end{align*}
where $(p',q',z')$ are the dual coordinates of
$(v_+,v_-,n_1)$ and $c$ is the $z$-exponent of the transition
monomial for $p'$ ($c=0$ for $(\PP^1)^3$, where the transition is
$(p,q,1/z)$, and $c=-1$ for $\Bl_{pt}\PP^3$, where it is
$(p/z,\,q/z,\,1/z)$).  The corner-chart cap is necessary: in the
$\Bl_{pt}\PP^3$ case the naive closure of
$\{|p|=\varepsilon,q=0,|z|\ge R\}$ collapses onto the torus-fixed
point of the corner chart, which \emph{lies on} $\Sb=\{p'=q'z'\}$;
the cap $T_3$ instead keeps $|p'-q'z'|=\varepsilon R^{c}>0$.  On
$T_2$, $t=-\varepsilon R^{\gamma}e^{i(\theta+\gamma\varphi)}$ has
constant modulus; the $r=0$ boundary is the boundary torus of $T_1$
and the $r=1$ boundary maps under the transition onto
$\partial T_3$, so with orientations of the caps chosen to cancel
$\partial T_2$, $T=T_1+T_2+T_3$ is a closed $3$-cycle in $M$.

The pullback of $\Omega_M$ to $T_1$ and $T_3$ vanishes identically
($dp\mapsto0$, respectively $dq'\mapsto0$, while $\Omega_M$ is
regular there), and on $T_2$ the pullback of
$z^{\gamma-1}dp\wedge dq\wedge dz/t^2$ is the \emph{constant} form
$dr\wedge d\theta\wedge d\varphi$: writing it out, the complex
Jacobian of the parametrization equals $t^2z^{1-\gamma}$ exactly.
Hence $\int_T\Omega_M=\int_0^1\!\int_0^{2\pi}\!\int_0^{2\pi}
dr\,d\theta\,d\varphi=4\pi^2$ up to the orientation sign.

Finally $b_3(M)=1$: from the pair sequence of $(V,M)$ and
$H^3(V)=0$, $b_3(M)$ is the corank of the restriction
$H^2(V;\QQ)\to H^2(\Sb;\QQ)$, which is one in both cases.
\end{proof}

\begin{lemma}\label{lem:D}
Let $p$ be a singular point of $X$ whose germ is the anticanonical
cone over $\dP_7$, or over $\dP_6$ with the induced local smoothing
through the one-parameter component.  Then
$H_3(M_p;\QQ)\;(\cong\QQ)\to H_3(X_s;\QQ)$ is injective for all small
$s\neq0$.
\end{lemma}

\begin{proof}
By Section~\ref{sec:miniversal} the induced arc lies in a single
reduced component $c$ of the miniversal base of the germ, off the
local discriminant for $s\neq0$, and by Section~\ref{sec:sweeps} the
Milnor fiber is the corresponding sweep complement, which for the two
channels in the statement is one of the complements of
Lemma~\ref{lem:D1}.  The component carries Pinkham's
$\CC^*$-equivariant swept family with parameters of weight one
(Section~\ref{sec:miniversal}: $T^1$ sits in the Gorenstein degree);
let $\Omega^{\mathrm{loc}}$ be an equivariant algebraic generator of
its relative dualizing sheaf, homogeneous of some integral weight
$a$.  To construct it, work in the graded local ring at the vertex.
The relative dualizing module is invertible, and graded Nakayama
selects from any algebraic local generator a homogeneous component
that still generates.  Its $\CC^*$-saturation contains the entire
fiber $b=1$: every point of that fiber tends to the vertex when the
positive-weight cone coordinates and the weight-one parameter are
scaled to zero.  Thus its restriction to $b=1$ is an algebraic
generator of $\omega_M$, hence a constant multiple of $\Omega_M$ by
the unit calculation following \eqref{eq:omegaM}; rescale it so that
$\Omega^{\mathrm{loc}}|_{b=1}=\Omega_M$.
For a horizontal family $\nu$ of classes in $H_3(M;\QQ)$,
\[
\int_{\nu(b)}\Omega^{\mathrm{loc}}_b\;=\;c(\nu)\cdot\chi(b),
\]
with $\chi$ the weight-$a$ monomial in the component parameter and
$c(\nu)=\int_\nu\Omega_M$: single-valuedness holds because the
$\CC^*$-action is a group action on the total space with parameter
weight one, so transport once around the base circle is the identity;
and $c(\cdot)$ is linear.  Writing
$\Omega=g\cdot\Omega^{\mathrm{loc}}$ near $p$ with $g$ a unit as in
Section~\ref{sec:relform}, let $b=b(s)$ be the actual parameter of the
component.  The cycle obtained from the fixed compact cycle of
Lemma~\ref{lem:D1} by the $\CC^*$-action lies in the local Milnor
piece and converges uniformly to $p$, because all cone coordinates
have positive weights.  Hence $g|_{\nu(s)}=g(p,0)+o(1)$ uniformly,
while homogeneity factors the same monomial $\chi(b(s))$ from the main
term and the error.  The global period is therefore
\[
\int_{\nu(s)}\Omega_s
=\chi(b(s))\bigl(c(\nu)\,g(p,0)+o(1)\bigr).
\]

If $\nu\neq0$ in $H_3(M;\QQ)$ had $\nu_*=0$ in $H_3(X_s;\QQ)$ for
arbitrarily small $s\neq0$, then a rational multiple of $\nu(s)$
would bound and the closed form $\Omega_s$ would have zero period,
while $\chi(b(s))\neq0$ off the discriminant; letting $s\to0$ along such
values gives $c(\nu)\,g(p,0)=0$ and, since $g(p,0)\neq0$,
$c(\nu)=0$, contradicting the injectivity of the period pairing in
Lemma~\ref{lem:D1}.  (That $H_3(M;\QQ)$ has rank one is what converts
this class-by-class exclusion into injectivity for all small $s$; a
higher-rank channel would require a uniform argument.)
\end{proof}

\section{Surgery membership}\label{sec:surgery}

\subsection{The decomposition}\label{sec:decomp}
Let $p_1,\dots,p_r$ be the singular points of $X$, with disjoint
closed cone neighborhoods $U_p$ chosen inside Milnor balls, and
$W=X\smallsetminus\bigcup_p\operatorname{int}U_p$, so
$\partial W=\bigsqcup_pL_p$.  Good representatives exist for
arbitrary isolated singularities
(Greuel~\cite[Definition~2.12, Lemma~2.13, Theorem~2.15]{Greuel19}),
and for small $s$ the family is smooth, hence $C^\infty$-trivial, over
a neighborhood of $W$; therefore
\[
X_s\ \cong\ W\ \cup_{\bigsqcup L_p}\ \bigsqcup_pM_p,
\qquad
\Xh\ \cong\ W\ \cup_{\bigsqcup L_p}\ \bigsqcup_pN_p,
\]
where $M_p$ is the Milnor fiber of the induced local smoothing and
$N_p$ a neighborhood of the exceptional set of the fixed crepant
resolution.  Both gluings are along the same links.  Only the
$C^\infty$-triviality over a neighborhood of $W$ is used here, so the
same decomposition of $X_s$ along $\bigsqcup_pL_p$ holds for an
arbitrary one-parameter deformation, with $X_s\cap U_p$ in place of
$M_p$ at any point whose germ is not smoothed; this is the form
Section~\ref{sec:partialmem} uses.

\begin{lemma}[membership]\label{lem:M}
For every $\gamma\in H_3(X_s;\QQ)$, the Mayer--Vietoris boundary
$\partial\gamma\in\bigoplus_pH_2(L_p;\QQ)$ lies in the space $K(X)$
of Theorem~\ref{thm:B-intro}: its components die in every
$H_2(M_p;\QQ)$, and their sum dies in $H_2(W;\QQ)$, hence in
$H_2(\Xh;\QQ)$.  Moreover, at a node $q$ the $L_q$-component of
$\partial\gamma$ is a fixed nonzero multiple of
$(\gamma\cdot\delta_q)$ times the generator, and at a cone point $p$
the boundary map identifies
$H_3(M_p,L_p;\QQ)\xrightarrow{\ \sim\ }R_p$.
\end{lemma}

\begin{proof}
The first statement is Mayer--Vietoris exactness for
$X_s=W\cup\bigsqcup M_p$, together with $W\subset\Xh$.  At a node,
$M\simeq T^*S^3$: $H_3(M,L;\QQ)\cong H^3(M;\QQ)=\QQ$ is generated by
a cotangent fiber $F$ with $\partial F$ the $S^2$-generator of
$H_2(S^2\times S^3)$ and $F\cdot\delta=1$; excision identifies the
$L_q$-component of $\partial\gamma$ with
$\partial$ of the image of $\gamma$ in $H_3(M_q,L_q)$, which equals
$(\gamma\cdot\delta_q)F$ by the perfect pairing.  At a cone point, exactness of the pair sequence makes the image of
$\partial\colon H_3(M,L;\QQ)\to H_2(L;\QQ)$ equal to
$\ker(H_2(L;\QQ)\to H_2(M;\QQ))=R_p$, while Lefschetz duality gives
$\dim H_3(M,L;\QQ)=b_3(M)$, which by Section~\ref{sec:sweeps} equals
$1$, $1$, $2$ in the three channels, the rank of $R_p$ computed in
Proposition~\ref{prop:kernels}.  A surjection between spaces of equal
finite dimension is an isomorphism, so $\partial$ is injective with
image $R_p$; equivalently $H_3(M)\to H_3(M,L)$ is zero, $H_3(L)\to
H_3(M)$ is surjective, and the skew intersection form of $H_3(M;\QQ)$
vanishes.
\end{proof}

\subsection{Boundary compatibility}\label{sec:boundary}
Lemma~\ref{lem:M} uses the subspaces
$R_p\subset H_2(L_p)$ through the \emph{sweep-side} description of
Proposition~\ref{prop:kernels} (the unit normal circle bundle of
$\Sb\subset V$), while the map to $H_2(\Xh)$ and all coordinates in
Section~\ref{sec:example} use the \emph{resolution-side} description
(the Seifert projection $L_p\to E_p$ of the cone link).  These are a
priori different identifications of $H_2(L_p;\QQ)$ with
$K^{\perp}$, and an arbitrary boundary diffeomorphism need not preserve
either marking.  Their compatibility is therefore a statement, not a
formality.  It appears not to be available in the literature:
the identification of the germ Milnor fiber with the sweep complement
is standard as \emph{open} manifolds
(\cite[Proposition~4.2]{PP15}), but the manifold-with-boundary form,
compatible with the canonical link identification, is what the
following lemma supplies.

\begin{lemma}[boundary compatibility]\label{lem:boundary}
Let $(X,p)$ be analytically isomorphic to the anticanonical cone over
the del Pezzo surface $S$, smoothed by a one-parameter family through
the component $c$ with sweep $V\supset\Sb$.  Then the boundary
identification $\partial M_p\cong L_p$ arising from the surgery
decomposition can be chosen so that the sweep-side projection
$H_2(L_p;\QQ)\to H_2(\Sb;\QQ)$ and the Seifert projection
$H_2(L_p;\QQ)\to H_2(E_p;\QQ)$ agree up to an isomorphism of
anticanonically polarized surfaces $(E_p,-K)\cong(\Sb,-K)$.  In
particular the comparison is a lattice isometry, and $R_p$, read in
resolution-side coordinates, is the root subspace of
Proposition~\ref{prop:kernels}.
\end{lemma}

\begin{proof}
Let $\mathcal Y_c\to B_c$ be the $\CC^*$-equivariant
Altmann--Pinkham family on the reduced component $c$, and let
$\overline{\mathcal Y}_c\to B_c$ be its projective sweep
compactification.  All component parameters have weight one, all
affine cone coordinates have positive weights, and the divisor at
infinity is the fixed family $S\times B_c$: negative weight does not
change the top graded piece.  Choose an equivariant affine embedding,
write $w_1,\dots,w_N>0$ for the coordinate weights and, for a common
multiple $d$ of them, use the weighted radius
\[
 \rho(z)=\left(\sum_{i=1}^N|z_i|^{2d/w_i}\right)^{1/(2d)},
 \qquad \rho(r\cdot z)=r\rho(z)\quad(r>0).
\]
Fix a sufficiently small weighted Milnor radius $\varepsilon$ and a
small sphere of radius $\eta$ in $B_c$.  Every sufficiently small
parameter $\beta\in c$ off the discriminant is uniquely
$\beta=t b$ with $t>0$ and $\lVert b\rVert=\eta$.  The discriminant is
homogeneous, so $b$ is also off it.  The action by $t^{-1}$ gives a
diffeomorphism of pairs
\begin{equation}\label{eq:compactcore}
 \left(\mathcal Y_\beta\cap\{\rho\leq\varepsilon\},
       \mathcal Y_\beta\cap\{\rho=\varepsilon\}\right)
 \ \cong\
 \left(K_{t,b},\Sigma_{t,b}\right),
 \quad
 K_{t,b}:=\mathcal Y_b\cap\{\rho\leq\varepsilon/t\}.
\end{equation}
The pair on the left is the actual compact Milnor fiber and its
boundary, not merely its interior.

We next identify the pair on the right with the sweep compact core.
Near the fixed divisor $S\subset\overline{\mathcal Y}_b$, choose a
normal coordinate $q$, and pass to the real oriented blowup along
$S$.  An affine coordinate of weight $w_i$ has the form
\[
 z_i=q^{-w_i}\bigl(a_i(y,b)+O(q)\bigr),
\]
where $y\in S$ and the leading coefficients do not vanish
simultaneously.  Consequently
\begin{equation}\label{eq:radiusatinfinity}
 \rho=|q|^{-1}A(y,\arg q,|q|,b)
\end{equation}
with $A$ smooth and strictly positive on the oriented boundary.  For
$t$ small, the equation $\rho=\varepsilon/t$ therefore has a unique
solution for $|q|$ and is a smooth radial graph over the unit normal
circle bundle of $S$.  Normal radial flow identifies this graph with
the boundary of a fixed tube of $S$, preserving the projection to
$S$, and extends across the intervening collar.  Thus
\eqref{eq:compactcore} continues to a diffeomorphism of pairs
\[
 (M_p,\partial M_p)\ \cong\
 \bigl(\overline{\mathcal Y}_b\smallsetminus
       \operatorname{int}\mathrm{Tube}(S),
       \partial\mathrm{Tube}(S)\bigr)
\]
that carries the actual boundary inclusion to the tube-boundary
inclusion.

By the sweep description of Section~\ref{sec:sweeps}, the compactified
smooth pair $(\overline{\mathcal Y}_b,S)$ is the corresponding
polarized sweep pair $(V,\Sb)$; along a connected smooth channel the
pair identifications may be transported by Ehresmann with the
trivializing vector field chosen to vanish on the fixed divisor $S$.

It remains to identify the marking of that boundary.  Apply the same
construction to the radial family $\mathcal Y_{\tau b}$ for
$0<\tau\leq t$.  Formula~\eqref{eq:radiusatinfinity} shows, on the
real oriented blowup, that the rescaled boundary graphs extend
smoothly to $\tau=0$.  Their limiting projection sends a nonzero point
of the central cone to its projective class in $S$, while the normal
circle is its $S^1$-orbit.  It is therefore exactly the Seifert
projection $L_p\to S$ of the cone link.  We choose the boundary
trivialization in the surgery decomposition to be this radial
transport.  The sweep-side tube projection and the resolution-side
Seifert projection then induce the same map on $H_2$, after the
polarized identification $(E_p,-K)\cong(S,-K)\cong(\Sb,-K)$.

Finally the $A_1$ and $A_2$ summands for $\dP_6$ and the unique
$(-2)$-line for $\dP_7$ are intrinsic to the polarized lattice, so
they correspond.  Since the displayed diffeomorphism is one of pairs,
it computes the kernel of the actual inclusion
$H_2(\partial M_p)\to H_2(M_p)$, as required.
\end{proof}

\subsection{Deformations that leave germs singular}\label{sec:partialmem}
Theorem~\ref{thm:B-intro} concerns smoothings.  A threefold carrying a
germ with no smoothing at all admits none, but it may still admit
deformations that smooth its other germs, and the surgery argument
constrains those as well.  We record the form of Lemma~\ref{lem:M}
that Section~\ref{sec:partner} uses.  Call an isolated singularity
\emph{reduced-rigid} if the reduction of its miniversal base is a
single point; this does not assert $T^1=0$.  The anticanonical cone
over $\mathbb F_1=\Bl_1\PP^2$, one of the two del Pezzo surfaces of
degree~$8$ (not $\PP^1\times\PP^1$, whose cone is smoothable), is
reduced-rigid: its quadrilateral, with edge vectors
$(1,0),(1,2),(-2,-1),(0,-1)$, admits no nontrivial Minkowski
decomposition \emph{into lattice polygons}, although its cone of
rational Minkowski summands is two-dimensional, so that $\dim T^1=1$
and the miniversal base is
$\operatorname{Spec}\CC[\varepsilon]/(\varepsilon^2)$
\cite[(9.1)(iv),\,(9.2)]{Altmann97}; see also
Gross~\cite[Example~2.19, \S5.5]{Gross97} and \cite{PaperI}.

\begin{lemma}\label{lem:rigid}
Let $(Y,0)$ be reduced-rigid and let $\cX\to\Disk$ be a
one-parameter deformation of a complex space $X$ with
$(X,p)\cong(Y,0)$.  Then the induced deformation of the germ at $p$ is
trivial.  Consequently there are a neighborhood $U$ of $p$ in $\cX$
and an $\varepsilon>0$ with $U\cong V\times\Disk_\varepsilon$ over
$\Disk_\varepsilon$, so that for $0<|s|<\varepsilon$ the fiber
$X_s\cap U$ has a singular point $q_s$, with $q_s\to p$ and
$(X_s,q_s)\cong(Y,0)$.
\end{lemma}

\begin{proof}
Near $p$ the family is induced from the miniversal deformation by a
holomorphic map $c\colon(\Disk,0)\to(B,0)$.  The analytic local ring
of a disk is a domain, hence reduced, so $c$ factors through
$B_{\mathrm{red}}$, a reduced point; thus $c$ is constant, and the
pullback of the miniversal family along a constant map is the product
$(Y,0)\times(\Disk,0)$.  The final statement is the passage from the
germ at $(p,0)$ to fibers, by choosing representatives of that
product decomposition.
\end{proof}

\begin{lemma}[membership, with germs left singular]\label{lem:Mpartial}
Let $X$ and $\Xh$ be as in the setting fixed in
Section~\ref{sec:intro} before Theorem~\ref{thm:B-intro}, except that
$X$ may carry in addition finitely many anticanonical cones over
$\mathbb F_1$, with an exceptional $\mathbb F_1$ over each of them on
$\Xh$; no smoothing of $X$ is assumed to exist.  Let $\cX\to\Disk$ be
any one-parameter deformation of $X$ and let
\[
\Sigma\ =\ \{\,p\in\Sing X\ :\ \text{the local fiber }X_s\cap U_p
\text{ is singular}\,\},
\]
a set independent of $s$ for small $s\neq0$, because each local
classifying arc either lies in the local discriminant or meets it
only at the origin.  Put
\[
X_s^{\flat}=W\cup\bigsqcup_{p\notin\Sigma}M_p ,
\]
a compact oriented $6$-manifold with boundary
$\bigsqcup_{p\in\Sigma}L_p$ contained in the smooth locus of $X_s$.
Then for every $\gamma\in H_3(X_s^{\flat},\partial X_s^{\flat};\QQ)$
there is a class $\kappa\in\bigoplus_pH_2(L_p;\QQ)$ with
\[
\kappa\in\ker\Bigl(\bigoplus_pH_2(L_p;\QQ)\to H_2(\Xh;\QQ)\Bigr),
\qquad
\kappa_p\in R_p\quad(p\notin\Sigma\text{ a cone point}),
\]
whose component at a node $q\notin\Sigma$ is a fixed nonzero multiple
of $(\gamma\cdot\delta_q)$ times the generator, and whose component at
a cone point $p\notin\Sigma$ is the image of $\gamma$ under the
composite
$H_3(X_s^{\flat},\partial X_s^{\flat};\QQ)\to H_3(M_p,L_p;\QQ)
\xrightarrow{\ \partial,\ \sim\ }R_p$.  Here $\gamma\cdot\delta_q$ is
the Poincar\'e--Lefschetz pairing of $X_s^{\flat}$.  No \emph{local}
condition is imposed on $\kappa_p$ for $p\in\Sigma$; those coordinates
are of course still coupled to the others by the displayed kernel
condition.
\end{lemma}

\begin{proof}
By Lemma~\ref{lem:rigid} every $\mathbb F_1$-cone point lies in
$\Sigma$, so a cone point $p\notin\Sigma$ is a $\dP_6$- or
$\dP_7$-cone point and the local computations of Lemma~\ref{lem:M}
apply to it.  Apply Mayer--Vietoris for the pair
$(X_s^{\flat},\partial X_s^{\flat})$ with $A=W$,
$B=\bigsqcup_{p\notin\Sigma}M_p$ and
$\partial X_s^{\flat}\subset A$, the couple being made excisive by the
standard collar thickening:
\[
H_3(X_s^{\flat},\partial X_s^{\flat};\QQ)
\xrightarrow{\ \partial\ }
\bigoplus_{p\notin\Sigma}H_2(L_p;\QQ)
\longrightarrow
H_2\bigl(W,\partial X_s^{\flat};\QQ\bigr)\oplus
\bigoplus_{p\notin\Sigma}H_2(M_p;\QQ).
\]
The second summand gives $(\partial\gamma)_p\in R_p$ at the cone
points $p\notin\Sigma$.  For the first, let $x\in H_2(W;\QQ)$ be the
image of $\partial\gamma$.  It dies in
$H_2(W,\partial X_s^{\flat};\QQ)$, so by exactness of
\[
H_2(\partial X_s^{\flat};\QQ)\longrightarrow H_2(W;\QQ)
\longrightarrow H_2(W,\partial X_s^{\flat};\QQ)
\]
it is the image of some
$y\in H_2(\partial X_s^{\flat};\QQ)=
\bigoplus_{p\in\Sigma}H_2(L_p;\QQ)$.  Set
$\kappa_p=(\partial\gamma)_p$ for $p\notin\Sigma$ and
$\kappa_p=-y_p$ for $p\in\Sigma$.  Then $\kappa$ dies in $H_2(W;\QQ)$,
hence in $H_2(\Xh;\QQ)$ because $W\subset\Xh$.

For the local statements, fix a node $q\notin\Sigma$.  Since
$\partial X_s^{\flat}\subset W$ is disjoint from $M_q$, excision gives
a restriction
$H_3(X_s^{\flat},\partial X_s^{\flat};\QQ)\to
H_3(M_q,L_q;\QQ)$, and under Poincar\'e--Lefschetz duality on
$X_s^{\flat}$ it is dual to the inclusion
$H_3(M_q;\QQ)\to H_3(X_s^{\flat};\QQ)$.  With $F\cdot\delta_q=1$ as in
the proof of Lemma~\ref{lem:M}, the restriction of $\gamma$ is
therefore $(\gamma\cdot\delta_q)F$, and $\kappa_q$ is
$(\gamma\cdot\delta_q)\,\partial F$.  The same excision at a cone
point $p\notin\Sigma$, followed by the isomorphism
$\partial\colon H_3(M_p,L_p;\QQ)\to R_p$ of Lemma~\ref{lem:M},
gives the second statement.
\end{proof}

\begin{remark}\label{rem:partialform}
The construction of $\Omega$ in Section~\ref{sec:relform} is unchanged
in this setting.  It uses the decomposition
\[
X_s\ =\ W\ \cup\ \bigsqcup_{p\notin\Sigma}M_p\ \cup\
\bigsqcup_{p\in\Sigma}\bigl(X_s\cap U_p\bigr)
\]
along all the links, the vanishing $H^1(L_p;\ZZ)=0$ of
Section~\ref{sec:relform}, which covers the $\mathbb F_1$ link, and the
triviality of $\omega_{\cX/\Disk}$ near each $(p,0)$.  The fiber $X_s$
is allowed to be singular.  That it is Gorenstein needs an argument
now that it is not smooth, and the following one replaces the appeal
to base change: $X_0$ is Gorenstein and is the divisor $\{t=0\}$ of
$\cX$ with $t$ a nonzerodivisor by flatness, so $\cX$ is Gorenstein
near $X_0$; each $X_s$ is then a Cartier divisor in a Gorenstein space
and hence Gorenstein, and adjunction makes $\omega_{\cX/\Disk}$
invertible with $\omega_{\cX/\Disk}|_{X_s}\cong\omega_{X_s}$.  The
singularities of $X_s$ are isolated for small $s$, since the singular
locus of $\cX\to\Disk$ is closed, proper over $\Disk$, and meets $X_0$
in a finite set; the same argument places $\Sing X_s$ inside
$\bigsqcup_p\operatorname{int}U_p$, which is what makes
$X_s^{\flat}$ lie in the smooth locus.  Since $X_s^{\flat}$ lies in
the smooth locus of $X_s$, the period arguments of Lemmas~\ref{lem:P}
and~\ref{lem:D} apply verbatim with $H_3(X_s;\QQ)$ replaced by
$H_3(X_s^{\flat};\QQ)$: requiring a bounding $4$-chain to lie in
$X_s^{\flat}$ is a stronger demand than requiring it in $X_s$, so the
conclusions are stronger, and the cycles involved lie in
$X_s^{\flat}$ by construction.  Poincar\'e--Lefschetz duality on the
compact oriented manifold with boundary $X_s^{\flat}$, which pairs
$H_3(X_s^{\flat};\QQ)$ perfectly with
$H_3(X_s^{\flat},\partial X_s^{\flat};\QQ)$, replaces Poincar\'e
duality on $X_s$.
\end{remark}

\section{The distinguished hypersurface and its finite data}
\label{sec:example}

\subsection{The threefold}\label{sec:X0}
Let $\Delta=\Delta_{\mathbb F_1}$ be the $22$-vertex reflexive
$4$-polytope of \cite[Theorem~6.1]{PaperIII} and
$\Delta^{\vee}$ its polar, with the $26$ vertices
$q_1,\dots,q_{26}$ listed in Table~\ref{tab:rays}.  The generic
anticanonical hypersurface
$X^{\circ}\subset\PP_{\Delta^{\vee}}$ has, for every
$\Delta$-regular choice of equation, exactly $26$ ordinary double
points (one for each unit-square $2$-face of $\Delta^\vee$), two
anticanonical $\dP_7$-cone points (the two pentagon faces with one
interior point) and two $\dP_6$-cone points (the two hexagon faces)
\cite[Section~6]{PaperIII}.  All thirty germs are \emph{analytically}
the cones: the ambient chart at the $3$-dimensional cone over a
$2$-face $F$ splits as $C(F)\times\CC^{*}$, the restriction of the
defining section to the one-dimensional orbit is
$t^{k}(c_a+c_bt)$ with $c_a,c_b\neq0$ by $\Delta$-regularity because
the dual edge $F^{\circ}$ has lattice length one, and Weierstrass
preparation in the orbit coordinate exhibits the germ as a graph over
$(C(F),0)$.  In particular each face contributes exactly one singular
point, the simple zero of $c_a+c_bt$.

$X^{\circ}$ satisfies the standing hypotheses of
Section~\ref{sec:periods}: it is an anticanonical Cartier divisor in
the Gorenstein Fano toric fourfold $P=\PP_{\Delta^{\vee}}$, so
$\omega_{X^{\circ}}\cong\cO_{X^{\circ}}$ by adjunction, and
$h^1(\cO_{X^{\circ}})=0$ from the sequence
$0\to\cO_P(K_P)\to\cO_P\to\cO_{X^{\circ}}\to0$ together with
$H^1(\cO_P)=0$ and
$H^2(\cO_P(K_P))\cong H^{2}(\cO_P)^{\vee}=0$ (toric vanishing and
Serre duality on the complete toric fourfold).

\begin{table}\centering
\caption{The $30$ lattice points of $\partial\Delta^{\vee}$: the
vertices $q_1,\dots,q_{26}$ and the four face-interior points:
$q_{27},q_{28}$ are the interior points of the two hexagons (below
$E_{6,1},E_{6,2}$) and $q_{29},q_{30}$ of the two pentagons (below
$E_{7,1},E_{7,2}$).}\label{tab:rays}
\begin{tabular}{llll}
\toprule
$q_1=(-1,-1,-1,-1)$ & $q_2=(-1,-1,-1,0)$ & $q_3=(-1,-1,0,-1)$ &
$q_4=(-1,-1,0,1)$\\
$q_5=(-1,-1,1,0)$ & $q_6=(-1,-1,1,1)$ & $q_7=(-1,0,-1,-1)$ &
$q_8=(-1,0,-1,0)$\\
$q_9=(-1,0,0,-1)$ & $q_{10}=(-1,0,0,1)$ & $q_{11}=(-1,0,1,0)$ &
$q_{12}=(-1,0,1,1)$\\
$q_{13}=(0,-1,-1,-1)$ & $q_{14}=(0,-1,-1,0)$ & $q_{15}=(0,-1,0,1)$ &
$q_{16}=(0,-1,0,0)$\\
$q_{17}=(0,1,-1,-1)$ & $q_{18}=(0,1,-1,0)$ & $q_{19}=(0,1,0,-1)$ &
$q_{20}=(0,1,0,0)$\\
$q_{21}=(1,1,-1,-1)$ & $q_{22}=(1,0,-1,0)$ & $q_{23}=(0,0,0,-1)$ &
$q_{24}=(0,0,0,1)$\\
$q_{25}=(0,0,1,0)$ & $q_{26}=(0,0,1,1)$ & $q_{27}=(-1,-1,0,0)$ &
$q_{28}=(-1,0,0,0)$\\
$q_{29}=(0,0,-1,-1)$ & $q_{30}=(0,0,-1,0)$ &&\\
\bottomrule
\end{tabular}
\end{table}

\subsection{The resolution and the matrix}\label{sec:matrix}
Fix the crepant resolution $\Xh\to X^{\circ}$ obtained from the
projective simplicial subdivision of the face fan of $\Delta^\vee$
using all thirty boundary lattice points, with the star subdivision at
the four face-interior points and the twenty-six square diagonals
\[
\begin{aligned}
&(q_2,q_7),(q_2,q_{13}),(q_3,q_7),(q_4,q_8),(q_4,q_{14}),
(q_5,q_9),(q_5,q_{13}),(q_5,q_{23}),(q_9,q_{23}),\\
&(q_6,q_{10}),(q_{10},q_{15}),(q_6,q_{11}),(q_6,q_{16}),
(q_6,q_{25}),(q_8,q_{17}),(q_9,q_{17}),(q_{12},q_{18}),
(q_{12},q_{24}),\\
&(q_{12},q_{19}),(q_{12},q_{25}),(q_{14},q_{16}),(q_{16},q_{23}),
(q_{16},q_{26}),(q_{18},q_{19}),(q_{20},q_{24}),(q_{20},q_{25}),
\end{aligned}
\]
labeled $N_1,\dots,N_{26}$ in this order; the subdivision is smooth
and projective (certified).  Over the nodes $\pi$ is small with
exceptional curves $\Gamma_i$, and over the cone points the
exceptional divisors are toric surfaces
$E_{6,1},E_{6,2}\cong\dP_6$ and $E_{7,1},E_{7,2}\cong\dP_7$.

We record the two families of functionals on $\Pic(Y_{\widehat\Sigma})$
(equivalently, by restriction, on $\Pic(\Xh)\otimes\QQ$, though only
one containment is used below) whose kernel computation is the finite
input.  \emph{Node rows}: for each square with cyclic vertices
$(a,b,c,d)$ and chosen diagonal $(a,c)$, the intersection functional
of $\Gamma_i$, with values $-1$ at $D_a,D_c$, $+1$ at $D_b,D_d$ and
$0$ otherwise; it annihilates principal divisors because
$q_a+q_c=q_b+q_d$.  \emph{Divisorial rows}: for each exceptional
surface, the composition of the restriction
$\Pic(Y)\to\Pic(E)$ ($D_j\mapsto$ the boundary curve of $E$ for the
rays of the face, $D_{\mathrm{int}}\mapsto K_E$, all other rays
$\mapsto0$) with a basis of the rank-$2$ or $3$ lattice
$K^{\perp}\subset\Pic(E)^{\vee}$.  In the bases of
Table~\ref{tab:kperp} the three $\dP_6$ functionals
$\varphi_{s_1},\varphi_{s_2},\varphi_{s_3}$ satisfy
$\varphi_{s_i}^2=-2$, $\varphi_{s_1}\cdot\varphi_{s_2}=1$,
$\varphi_{s_3}\perp\varphi_{s_1},\varphi_{s_2}$, so
$\langle\varphi_{s_3}\rangle$ is the $A_1$ summand and
$\langle\varphi_{s_1},\varphi_{s_2}\rangle$ the $A_2$ summand; the
two $\dP_7$ functionals satisfy $\varphi_\beta^2=-2$,
$\varphi_\alpha^2=-4$, and $\varphi_\beta$ spans the unique
$(-2)$-line.  Stacking the $26$ node rows and the ten divisorial rows
gives an integral matrix
$B\in\operatorname{Mat}_{36\times26}(\ZZ)$ in the dual basis of the
$26$ divisor classes complementary to the unimodular cone
$(q_{15},q_{22},q_{24},q_{26})$.

\begin{table}\centering
\caption{Exceptional-surface data.  Bases of $\Pic(E)$: boundary
divisors at the rays listed; intersection matrices $I$; canonical
classes $K$; and the $K^{\perp}$ functionals in the dual bases.  The
two $\dP_6$ surfaces carry identical data at rays
$(q_3,q_4,q_5,q_6)$, resp.\ $(q_9,q_{10},q_{11},q_{12})$.}
\label{tab:kperp}
\begin{tabular}{lll}
\toprule
surface & lattice data & $K^\perp$ rows\\
\midrule
$E_{6,i}$ &
$I_6=\begin{psmallmatrix*}[r]-1&0&1&0\\0&-1&0&1\\1&0&-1&1\\0&1&1&-1
\end{psmallmatrix*}$,\quad $K=(-1,-1,-2,-2)$ &
$\begin{aligned}\varphi_{s_1}&=(0,0,-1,1)\\
\varphi_{s_2}&=(-1,-1,1,0)\\
\varphi_{s_3}&=(-1,1,1,-1)\end{aligned}$\\[1.4em]
$E_{7,1}$ &
$I_{7,1}=\begin{psmallmatrix*}[r]0&0&1\\0&-1&1\\1&1&0
\end{psmallmatrix*}$ at $(q_{13},q_{17},q_{21})$,\quad
$K=(-1,-1,-2)$ &
$\begin{aligned}\varphi_{\alpha}&=(1,-1,0)\\
\varphi_{\beta}&=(1,1,-1)\end{aligned}$\\[1.2em]
$E_{7,2}$ &
$I_{7,2}=\begin{psmallmatrix*}[r]-1&0&1\\0&0&1\\1&1&0
\end{psmallmatrix*}$ at $(q_{14},q_{18},q_{22})$,\quad
$K=(-1,-1,-2)$ &
$\begin{aligned}\varphi_{\alpha}&=(-2,0,1)\\
\varphi_{\beta}&=(-1,-1,1)\end{aligned}$\\
\bottomrule
\end{tabular}
\end{table}

\subsection{The certified kernel computation}\label{sec:kernelcomp}
Write coordinates on $\QQ^{36}$ as
\[
(u_1,\dots,u_{26};\ s_{1},s_{2},s_{3}\text{ at }E_{6,1};\
s_{1},s_{2},s_{3}\text{ at }E_{6,2};\ \alpha_1,\beta_1;\
\alpha_2,\beta_2).
\]
The following statements are exact integral
computations, reproduced independently by the ancillary programs from
the printed data of Tables~\ref{tab:rays} and~\ref{tab:kperp} together
with the subdivision data document (see
Appendix~\ref{app:cert}):
\begin{enumerate}
\item $\rk B=21$ and $\dim\ker(B^{\mathsf T})=15$; the node
subsystem has rank $18$, with coloops exactly $N_9$ and $N_{11}$; the
unique row coloop of the full matrix is the $\beta_1$-row, i.e.\
$\beta_1$ vanishes identically on $\ker(B^{\mathsf T})$.
\item On the \emph{enlarged restricted kernel}
$\ker(B^{\mathsf T})\cap\{\alpha_1=\alpha_2=0\}$, of dimension $13$,
the four coordinates $u_9,u_{11},\beta_1,\beta_2$ vanish identically.
The four finer profiles obtained by additionally restricting each
$\dP_6$ block to its $A_1$ line ($s_1=s_2=0$) or its $A_2$ plane
($s_3=0$) have dimensions $9,10,10,12$ and force the same four
coordinates (and one $s_3$ on each mixed profile).
\item The forcing is critical in the $\dP_7$ data: releasing either
$\alpha$-coordinate destroys it.
\end{enumerate}
By construction $\ker(B^{\mathsf T})$ contains the kernel of
$\bigoplus_pH_2(L_p;\QQ)\to H_2(\Xh;\QQ)$ written in these
coordinates: the node generators map to $\pm[\Gamma_i]$ and the
divisorial classes to their pushforwards from
$K^{\perp}\subset H_2(E_p)$, and each row is the pairing with an
ambient divisor class.  Only this containment is used: if the ambient
Picard group under-detected $H_2(\Xh;\QQ)$, the certified kernel
would be larger than the topological one and every forced-zero
statement would still bind.  (In fact equality holds:
$(-K_Y)^4=94\neq0$ and $h^{1,1}(\Xh)=26$; but the proofs below do not
consume this.)

\section{Proofs of the main results}\label{sec:proofs}

\begin{theorem}[mixed necessity]\label{thm:necessity}
Let $X$ be a compact complex threefold with $\omega_X\cong\cO_X$ and
$h^1(\cO_X)=0$ whose singular points are ordinary double points and
anticanonical cones over $\dP_6$ or $\dP_7$, let
$\Xh\to X$ be a crepant resolution as in
Section~\ref{sec:intro}, and suppose $X$ admits a one-parameter
smoothing.  With $R_p=\ker(H_2(L_p;\QQ)\to H_2(M_p;\QQ))$ for the
actual Milnor fibers of the induced local smoothings, set
\[
K(X)=\Bigl\{\kappa\in\ker\Bigl(\bigoplus_pH_2(L_p;\QQ)
\to H_2(\Xh;\QQ)\Bigr)\ :\ \kappa_p\in R_p
\text{ at every cone point }p\Bigr\}.
\]
Then:
\begin{enumerate}
\item for every node $q$ there is $\kappa\in K(X)$ with
$\kappa_q\neq0$;
\item for every $\dP_7$ point $q$ there is $\kappa\in K(X)$ whose
$q$-component is a nonzero element of the $(-2)$-line $R_q$;
\item for every $\dP_6$ point $q$ smoothed through the one-parameter
component of its miniversal base there is $\kappa\in K(X)$ whose
$q$-component is a nonzero element of $R_q$ (the $A_1$ summand).
\end{enumerate}
\end{theorem}

\begin{proof}
Fix small $s\neq0$.  (1)~By Lemma~\ref{lem:P}, $\delta_q\neq0$ in
$H_3(X_s;\QQ)$; by Poincar\'e duality on the closed oriented $X_s$
choose $\gamma$ with $\gamma\cdot\delta_q\neq0$; by
Lemma~\ref{lem:M}, $\kappa=\partial\gamma$ lies in $K(X)$ and its
$q$-component is a nonzero multiple of $\gamma\cdot\delta_q$.
(2),(3)~By Lemma~\ref{lem:D} the generator $\nu$ of
$H_3(M_q;\QQ)$ survives in $H_3(X_s;\QQ)$; choose $\gamma$ with
$\gamma\cdot\nu_*\neq0$; by Lemma~\ref{lem:M} the $q$-component of
$\partial\gamma$ is the image of a nonzero element of
$H_3(M_q,L_q;\QQ)$ under the isomorphism onto $R_q$, and $R_q$ is
identified with the stated root subspace by
Proposition~\ref{prop:kernels}, Section~\ref{sec:match} and
Lemma~\ref{lem:boundary}.
\end{proof}

\begin{remark}\label{rem:single}
Since a $\QQ$-vector space is not the union of finitely many proper
subspaces, the per-point conclusions of
Theorem~\ref{thm:necessity} combine: there is a \emph{single} class
$\kappa\in K(X)$ simultaneously nonzero at every node, every
$\dP_7$ point, and every line-smoothed $\dP_6$ point.
\end{remark}

\begin{corollary}\label{cor:friedman}
In the purely nodal case, Theorem~\ref{thm:necessity}(1) with
Remark~\ref{rem:single} states: any one-parameter smoothing forces a
relation among the classes $[\Gamma_i]\in H_2(\Xh;\QQ)$ of the
exceptional curves of a small resolution in which every node appears
with nonzero coefficient.  \qed
\end{corollary}

\begin{proposition}\label{prop:coloops}
If $X^{\circ}$ admitted a one-parameter smoothing, the vanishing
spheres of the nodes $N_9$ and $N_{11}$ would satisfy
$\delta_9=\delta_{11}=0$ in $H_3(X_s;\QQ)$.
\end{proposition}

\begin{proof}
By Lemma~\ref{lem:M} every boundary class
$\partial\gamma$ lies in $K(X^{\circ})$; its $\dP_7$-components lie
on the $(-2)$-lines by Proposition~\ref{prop:kernels}, and
Lemma~\ref{lem:boundary} identifies these with the $(-2)$-lines of
the resolution-side coordinates of Section~\ref{sec:kernelcomp}, so
$\partial\gamma$ lies in the enlarged restricted kernel, on which
$u_9$ and $u_{11}$ vanish identically (item~(2) there).  Since the $u_q$-component of $\partial\gamma$ is a
nonzero multiple of $\gamma\cdot\delta_q$, nondegeneracy of the
intersection pairing gives $\delta_9=\delta_{11}=0$.
\end{proof}

\begin{theorem}\label{thm:main}
$X^{\circ}$ admits no one-parameter smoothing.
\end{theorem}

\begin{proof}
Suppose one exists.  Every germ of $X^{\circ}$ is then smoothed; in
particular the node $N_9$ is, so Lemma~\ref{lem:P} gives
$\delta_9\neq0$ in $H_3(X_s;\QQ)$, contradicting
Proposition~\ref{prop:coloops}.

Alternatively, avoiding the nodes altogether: the point covered by
$E_{7,1}$ is smoothed through the unique component of its miniversal
base, so Theorem~\ref{thm:necessity}(2) produces
$\kappa\in K(X^{\circ})$ whose $E_{7,1}$-component is a nonzero
element of the $(-2)$-line, i.e.\ (Lemma~\ref{lem:boundary} and
Table~\ref{tab:kperp}) has nonzero $\beta_1$-coordinate; but
$\beta_1$ vanishes identically already on the unrestricted
$\ker(B^{\mathsf T})$ (Section~\ref{sec:kernelcomp}, item~(1)).
\end{proof}

\begin{remark}\label{rem:independence}
The two proofs are independent in the period lemma and in the
certified fact consumed (Lemma~\ref{lem:P} with
Section~\ref{sec:kernelcomp}(2), versus Lemma~\ref{lem:D} with
Section~\ref{sec:kernelcomp}(1)), and share the membership and
boundary-compatibility inputs (Lemmas~\ref{lem:M}
and~\ref{lem:boundary}) and the divisorial rows of
Table~\ref{tab:kperp}.
\end{remark}

\begin{corollary}\label{cor:def}
No fiber of a small representative of the miniversal deformation of
$X^{\circ}$ is smooth; consequently $X^{\circ}$ is not smoothable over
any base germ, and neither is any small deformation of $X^{\circ}$.
Likewise, no formal arc in the miniversal deformation of
$X^{\circ}$ meets all thirty local smoothing loci.
\end{corollary}

\begin{proof}
The miniversal deformation of the compact $X^{\circ}$ exists
\cite{Grauert74} and its discriminant is a closed analytic subset of
the base.  If some fiber arbitrarily close to the origin were smooth,
the curve selection lemma \cite{Milnor68} would produce a
real-analytic arc through the origin with smooth generic fiber;
complexifying (a real-analytic function vanishing on a real segment
of a disk vanishes identically, and the discriminant preimage of the
complexified arc is discrete) gives a holomorphic disk, that is, a
one-parameter smoothing, contradicting Theorem~\ref{thm:main}.  Small
deformations are covered by openness of versality.  For the formal
statement: by Artin approximation \cite{Artin68} applied to the
analytic miniversal deformation, a formal arc whose local components
avoid all thirty local discriminants would be approximated, to any
jet order beyond the vanishing orders of those components, by an
analytic arc with the same property.  Since the singular locus of the
fibers is closed in the total space and meets the central fiber only
in the thirty points, the fibers of the approximating arc are, for
small $t$, singular at most inside the thirty Milnor-ball
neighborhoods; avoidance of the local discriminants makes them smooth
there, hence smooth everywhere, so the arc is a one-parameter
smoothing, contradicting Theorem~\ref{thm:main}.  (The local
components are read
through choices of local classifying maps; avoidance of the local
discriminants is choice-independent, being smoothness of the local
fibers.)
\end{proof}

\begin{remark}\label{rem:mirror}
The two nodal coloops quantify how far $X^{\circ}$ is from the nodal
theory: no relation among the twenty-six exceptional curves alone is
nonzero at $N_9$ or $N_{11}$, so a smoothing would have had to route
through the divisorial directions, and items~(2)--(3) of
Section~\ref{sec:kernelcomp} say the divisorial directions cannot
supply what is needed at any order.  The mirror partner of
$X^{\circ}$ is treated in Section~\ref{sec:partner}.
\end{remark}

\section{The mirror partner}\label{sec:partner}
The polytope $\Delta$ itself carries a Calabi--Yau hypersurface, the
mirror partner of $X^{\circ}$ in the doubly isolated pair.  Write
$X_{\Delta}\subset\PP_{\Delta}$ for a $\Delta$-regular anticanonical
hypersurface and $v_1,\dots,v_{22}$ for the vertices of $\Delta$ in
the order of \cite[Theorem~6.1]{PaperIII}, reproduced in
Table~\ref{tab:partnerrays}.  By that theorem $X_{\Delta}$ has twenty
ordinary double points, one $\dP_6$-cone point over the hexagonal face
$\conv\{v_9,v_{14},v_{15},v_{20},v_{21},v_{22}\}$, and one point whose
germ is the anticanonical cone over $\mathbb F_1$, over the
quadrilateral face $\conv\{v_4,v_{18},v_{19},v_{22}\}$.  As in
Section~\ref{sec:X0}, $X_{\Delta}$ is an anticanonical Cartier divisor
in a Gorenstein Fano toric fourfold, so
$\omega_{X_{\Delta}}\cong\cO_{X_{\Delta}}$ and
$h^1(\cO_{X_{\Delta}})=0$, and every germ is analytically the cone
over its face because every dual edge has lattice length one.  The
last germ
is reduced-rigid, so $X_{\Delta}$ admits no smoothing \cite{PaperI}.
Theorem~\ref{thm:main} and the reduced-rigidity of this germ together
suggest a division of labour, $X_{\Delta}$ failing to smooth locally
and $X^{\circ}$ globally.  The division does not hold up: the
obstruction of this paper is present on $X_{\Delta}$ as well, in the form
Theorem~\ref{thm:partial} gives it, and it does not use the
$\mathbb F_1$ germ.

\begin{theorem}[necessity with germs left singular]\label{thm:partial}
Let $X$ be a compact complex threefold with $\omega_X\cong\cO_X$ and
$h^1(\cO_X)=0$ whose singular points are ordinary double points,
anticanonical cones over $\dP_6$ or $\dP_7$, and anticanonical cones
over $\mathbb F_1$; let $\Xh\to X$ be a crepant resolution as in
Section~\ref{sec:intro}.  Let $\cX\to\Disk$ be any one-parameter
deformation of $X$ and let $\Sigma$ be the set of points at which the
nearby fiber is singular, a set containing every $\mathbb F_1$-cone
point by Lemma~\ref{lem:rigid}.  Put
\[
K^{\flat}(X;\Sigma)=
\Bigl\{\kappa\in\ker\Bigl(\bigoplus_pH_2(L_p;\QQ)\to H_2(\Xh;\QQ)
\Bigr)\ :\ \kappa_p\in R_p\text{ at every cone point }p\notin\Sigma
\Bigr\}.
\]
Then $K^{\flat}(X;\Sigma)$ contains a class nonzero at $q$ for every
node $q\notin\Sigma$, for every $\dP_7$ point $q\notin\Sigma$, and for
every $\dP_6$ point $q\notin\Sigma$ smoothed through the one-parameter
component of its miniversal base.  For $\Sigma=\varnothing$ this is
Theorem~\ref{thm:necessity}.
\end{theorem}

\begin{proof}
The proof of Theorem~\ref{thm:necessity} applies with
Remark~\ref{rem:partialform} and Lemma~\ref{lem:Mpartial} in place of
Lemma~\ref{lem:M}.  At a node $q\notin\Sigma$, Lemma~\ref{lem:P}
gives $\delta_q\neq0$ in $H_3(X_s^{\flat};\QQ)$, since a bounding
$4$-chain in $X_s^{\flat}$ would make the period of the closed form
$\Omega_s$ vanish; Poincar\'e--Lefschetz duality on the compact
oriented manifold with boundary $X_s^{\flat}$ produces
$\gamma\in H_3(X_s^{\flat},\partial X_s^{\flat};\QQ)$ with
$\gamma\cdot\delta_q\neq0$; and Lemma~\ref{lem:Mpartial} converts
$\gamma$ into a class $\kappa\in K^{\flat}(X;\Sigma)$ with
$\kappa_q\neq0$.  At the cone points $q\notin\Sigma$ in the two
covered channels the same substitution is made in the argument of
Theorem~\ref{thm:necessity}(2),(3), with Lemma~\ref{lem:D} read in
$H_3(X_s^{\flat};\QQ)$.
\end{proof}

Whatever the deformation and whatever $\Sigma$, the space
$K^{\flat}(X;\Sigma)$ is by definition contained in
$\ker\bigl(\bigoplus_pH_2(L_p;\QQ)\to H_2(\Xh;\QQ)\bigr)$, so a
coordinate that vanishes identically on the latter is forced in every
case.  On $X_{\Delta}$ that is what happens.

Fix the smooth projective crepant subdivision of the face fan of
$\Delta$ that uses all twenty-four boundary lattice points, with the
star subdivision at $v_{23}$ and $v_{24}$ and the twenty square
diagonals
\[
\begin{aligned}
&(v_2,v_7),(v_2,v_9),(v_3,v_8),(v_3,v_9),(v_{10},v_{12}),
(v_{10},v_{15}),(v_{11},v_{13}),(v_{11},v_{14}),(v_{11},v_{17}),\\
&(v_{10},v_{16}),(v_9,v_{12}),(v_{12},v_{17}),(v_{14},v_{17}),
(v_9,v_{13}),(v_{13},v_{16}),(v_{15},v_{16}),(v_{15},v_{19}),
(v_{14},v_{18}),\\
&(v_{18},v_{20}),(v_{19},v_{21}),
\end{aligned}
\]
labeled $n_1,\dots,n_{20}$; the subdivision is smooth, with $94$
maximal cones, and projective (certified).  The corresponding circuit
relations are
\[
\begin{aligned}
&v_1{+}v_{12}=v_2{+}v_7,\ \
v_1{+}v_{15}=v_2{+}v_9,\ \
v_1{+}v_{13}=v_3{+}v_8,\ \
v_1{+}v_{14}=v_3{+}v_9,\\
&v_2{+}v_{19}=v_{10}{+}v_{12},\ \
v_2{+}v_{20}=v_{10}{+}v_{15},\ \
v_3{+}v_{18}=v_{11}{+}v_{13},\ \
v_3{+}v_{21}=v_{11}{+}v_{14},\\
&v_5{+}v_{21}=v_{11}{+}v_{17},\ \
v_6{+}v_{20}=v_{10}{+}v_{16},\ \
v_7{+}v_{15}=v_9{+}v_{12},\ \
v_7{+}v_{19}=v_{12}{+}v_{17},\\
&v_7{+}v_{21}=v_{14}{+}v_{17},\ \
v_8{+}v_{14}=v_9{+}v_{13},\ \
v_8{+}v_{18}=v_{13}{+}v_{16},\ \
v_8{+}v_{20}=v_{15}{+}v_{16},\\
&v_{12}{+}v_{20}=v_{15}{+}v_{19},\ \
v_{13}{+}v_{21}=v_{14}{+}v_{18},\ \
v_{16}{+}v_{22}=v_{18}{+}v_{20},\ \
v_{17}{+}v_{22}=v_{19}{+}v_{21},
\end{aligned}
\]
in the order $n_1,\dots,n_{20}$; in each of them the right-hand pair
is the chosen diagonal, so the circuit row of
Section~\ref{sec:matrix} is $-1$ there and $+1$ on the left-hand pair.
Choosing the other diagonal negates the row and changes nothing
below.
Table~\ref{tab:partnersurf} records the two exceptional surfaces.  The
lattice $K^{\perp}(\mathbb F_1)$ is a line spanned by a primitive
class of square $-8$, and it contains no $(-2)$-class, consistently
with the reduced-rigidity of that germ.  No condition is imposed at
that point in any case: by Lemma~\ref{lem:rigid} it lies in $\Sigma$
for every deformation, so no Milnor fiber and no $R_p$ is defined
there.  Stacking the twenty circuit rows, the three hexagon rows and
the single $\mathbb F_1$ row gives an integral
matrix $B_{\Delta}\in\operatorname{Mat}_{24\times24}(\ZZ)$, the
columns being indexed by the divisors of the twenty-four boundary
lattice points, which span $\Pic$ of the resolution without being a
basis.  The
following are exact integral computations, reproduced by the ancillary
programs from the printed data of
Tables~\ref{tab:partnerrays} and~\ref{tab:partnersurf}:
\begin{enumerate}
\item $\rk B_{\Delta}=16$ and $\dim\ker(B_{\Delta}^{\mathsf T})=8$;
the node subsystem has rank $14$, with coloops exactly $n_9$ and
$n_{10}$;
\item on the unrestricted kernel $\ker(B_{\Delta}^{\mathsf T})$ the
coordinates of $n_9$, of $n_{10}$, of the $\mathbb F_1$ row and of the
$A_1$ row of the hexagon vanish identically, so no branch restriction
is needed;
\item the column of $B_{\Delta}$ at the divisor of $v_{24}$ is
identically zero, the node rows having no entry there and the
divisorial rows having entry $\langle\varphi,K_E\rangle=0$; the same
holds at $v_{23}$.  Discarding the $\mathbb F_1$ row and that column
altogether leaves a system on which $n_9$ and $n_{10}$ are still
forced.
\end{enumerate}
As in Section~\ref{sec:kernelcomp}, only the containment of
$\ker(\bigoplus_pH_2(L_p;\QQ)\to H_2(\Xh;\QQ))$ in
$\ker(B_{\Delta}^{\mathsf T})$ is used.

\begin{theorem}\label{thm:partner}
No one-parameter deformation of $X_{\Delta}$ has nearby fiber smooth
at the node $n_9=\conv\{v_5,v_{11},v_{17},v_{21}\}$, and none has
nearby fiber smooth at $n_{10}=\conv\{v_6,v_{10},v_{16},v_{20}\}$.
The same two conclusions hold for deformations over arbitrary analytic
base germs.  In particular no analytic deformation of $X_{\Delta}$
smooths all twenty nodes and the $\dP_6$ point, whatever it does at the
$\mathbb F_1$ point.  By item~(2) and
Theorem~\ref{thm:partial}, the same argument applied to the $A_1$
coordinate shows that no analytic deformation of $X_{\Delta}$ smooths
its $\dP_6$ point through the one-parameter component of its miniversal
base.  The reduced-rigidity of the $\mathbb F_1$ germ excludes the
smoothings of $X_{\Delta}$; the deformations excluded here are ones
that it leaves available.
\end{theorem}

\begin{proof}
Suppose the nearby fiber is smooth at $n_9$, so $n_9\notin\Sigma$.  By
Theorem~\ref{thm:partial} there is
$\kappa\in K^{\flat}(X_{\Delta};\Sigma)$ with $\kappa_{n_9}\neq0$.
But $K^{\flat}(X_{\Delta};\Sigma)$ is contained in
$\ker(B_{\Delta}^{\mathsf T})$ for every $\Sigma$, read in the
coordinates above, and the $n_9$ coordinate vanishes identically
there.  The argument for $n_{10}$ is the same.

For the $\dP_6$ assertion, suppose first that the base is a disk and
that the local classifying arc lies in the one-parameter component of
the miniversal base.  Theorem~\ref{thm:partial} then gives a class whose
$A_1$ coordinate is nonzero, contrary to item~(2).

Now let the deformation have an arbitrary analytic base germ $(S,0)$.
Since the assertions concern actual nearby fibers, we may replace $S$
by $S_{\mathrm{red}}$.  If one of the excluded local smoothings occurred
arbitrarily close to the origin, the curve-selection argument used in
Corollary~\ref{cor:def}, applied to the complement of the inverse image
of the relevant local discriminant, would give a holomorphic disk
through the origin whose punctured fibers have the same local smoothing
property.  For the $\dP_6$ assertion, apply curve selection inside the
inverse image of the one-parameter component and outside its
discriminant.  Pullback to the resulting disk contradicts the
one-parameter conclusions just proved.
\end{proof}

\begin{corollary}\label{cor:pair}
Both members of the unique doubly isolated mixed mirror pair carry the
obstruction of this paper: on $X^{\circ}$, in the form of
Theorem~\ref{thm:necessity}, it excludes every smoothing; on
$X_{\Delta}$, in the form of Theorem~\ref{thm:partial}, it excludes
every analytic deformation smoothing $n_9$ or $n_{10}$.  On
$X_{\Delta}$ the two mechanisms are independent: by item~(3) the
forcing is computed from the twenty nodes and the hexagon alone, with the
$\mathbb F_1$ data discarded, whereas the reduced-rigidity of that germ
is a statement about it alone.
\end{corollary}

\begin{table}\centering
\caption{The $24$ boundary lattice points of $\Delta=\Delta_{\mathbb
F_1}$: the vertices $v_1,\dots,v_{22}$ in the order of
\cite[Theorem~6.1]{PaperIII}, the interior point $v_{23}$ of the
hexagonal face and the interior point $v_{24}$ of the $\mathbb F_1$
face.}\label{tab:partnerrays}
\begin{tabular}{llll}
\toprule
$v_1=(1,0,0,0)$ & $v_2=(0,1,0,0)$ & $v_3=(1,-1,0,0)$ &
$v_4=(0,0,1,0)$\\
$v_5=(0,0,0,1)$ & $v_6=(0,0,1,-1)$ & $v_7=(0,0,-1,1)$ &
$v_8=(0,0,0,-1)$\\
$v_9=(0,0,-1,0)$ & $v_{10}=(-1,1,0,0)$ & $v_{11}=(0,-1,0,0)$ &
$v_{12}=(-1,1,-1,1)$\\
$v_{13}=(0,-1,0,-1)$ & $v_{14}=(0,-1,-1,0)$ & $v_{15}=(-1,1,-1,0)$ &
$v_{16}=(-1,0,0,-1)$\\
$v_{17}=(-1,0,-1,1)$ & $v_{18}=(-1,-1,0,-1)$ & $v_{19}=(-2,1,-1,1)$ &
$v_{20}=(-2,1,-1,0)$\\
$v_{21}=(-1,-1,-1,0)$ & $v_{22}=(-2,0,-1,0)$ & $v_{23}=(-1,0,-1,0)$ &
$v_{24}=(-1,0,0,0)$\\
\bottomrule
\end{tabular}
\end{table}

\begin{table}\centering
\caption{The two exceptional surfaces of $\Xh_{\Delta}$: the cyclic
boundary rays, the intersection matrix $I$ in the Picard basis of the
boundary divisors at the rays listed, the canonical class $K$, and the
$K^{\perp}$ functionals in the dual basis.  As in Table~\ref{tab:kperp},
$\varphi_{s_3}$ spans the $A_1$ summand and
$\varphi_{s_1},\varphi_{s_2}$ the $A_2$ summand, here with
$\varphi_{s_1}\cdot\varphi_{s_2}=-1$.  The $\mathbb F_1$ lattice
$K^{\perp}$ has no $(-2)$-class.}\label{tab:partnersurf}
\begin{tabular}{lll}
\toprule
surface & lattice data & $K^{\perp}$ rows\\
\midrule
$E_{6,\Delta}$ over $v_{23}$, rays &
$I_6=\begin{psmallmatrix*}[r]-1&1&0&0\\1&-1&1&0\\0&1&-1&1\\
0&0&1&-1\end{psmallmatrix*}$ at $(v_9,v_{15},v_{20},v_{22})$, &
$\begin{aligned}\varphi_{s_1}&=(1,-1,0,1)\\
\varphi_{s_2}&=(0,-1,1,0)\\
\varphi_{s_3}&=(1,-1,1,-1)\end{aligned}$\\
$(v_{21},v_{14},v_9,v_{15},v_{20},v_{22})$ &
$K=(-1,-2,-2,-1)$ &\\[1.2em]
$E_{\mathbb F_1}$ over $v_{24}$, rays &
$I_{\mathbb F_1}=\begin{psmallmatrix*}[r]0&1\\1&-1
\end{psmallmatrix*}$ at $(v_{19},v_{22})$, &
$\varphi_{t}=(-2,3)$\\
$(v_{18},v_4,v_{19},v_{22})$ & $K=(-3,-2)$ &\\
\bottomrule
\end{tabular}
\end{table}

\section{The criterion as a finite test, and an open problem}
\label{sec:test}

Theorem~\ref{thm:necessity} is a finite test for any concrete
candidate: compute the matrix of Section~\ref{sec:matrix} from the
polytope (node circuit rows; $K^{\perp}$ bases of the star fans of
the pentagon and hexagon faces), restrict to the branch subspaces,
and check whether some covered coordinate is forced to vanish.  A
forced node coordinate contradicts Lemma~\ref{lem:P}; a forced
$\dP_7$ root coordinate contradicts Lemma~\ref{lem:D}; a forced
$s_3$-coordinate contradicts Lemma~\ref{lem:D} on profiles through
the line component.  The test consumes nothing about the example
beyond $\Delta$-regularity and its face inventory.

The test fires on threefolds that no purely nodal argument can reach.
Let $\Delta_{20}$ be the $20$-vertex reflexive polytope
\[
\begin{gathered}
\conv\{(1,0,0,0),(0,1,0,0),(0,0,1,0),(0,0,0,1),(1,-1,-1,-1),
(-1,1,1,1),(0,0,0,-1),\\
(0,0,-1,0),(0,-1,0,0),(-1,1,0,1),(0,0,-1,-1),(-1,1,0,0),
(-1,0,1,1),(0,-1,0,-1),\\
(-1,0,1,0),(0,-1,-1,-1),(-1,0,0,1),(-1,0,-1,0),(-1,-1,0,0),
(-1,-1,1,0)\},
\end{gathered}
\]
one of the both-sides-unit polytopes of \cite{PaperIII}.  Its generic
anticanonical hypersurface $X_{20}$ has seventeen ordinary double
points and one $\dP_7$-cone point, all isolated; as in
Section~\ref{sec:X0} it satisfies
$\omega_{X_{20}}\cong\cO_{X_{20}}$ and $h^1(\cO_{X_{20}})=0$, and a
smooth projective crepant subdivision exists, obtained as in
Section~\ref{sec:matrix} from the seventeen square diagonals and the
star subdivision at the interior point of the pentagon.  Here
$\Delta_{20}$ has $22$ lattice points and no facet with an interior
point, so Batyrev's formulas give
$h^{1,1}(\Xh_{20})=22-5=17=\operatorname{rank}\Pic$ of the ambient
toric fourfold: the ambient divisor classes span
$H^2(\Xh_{20};\QQ)$, and the certified kernel of
Section~\ref{sec:kernelcomp} coincides with the topological one rather
than merely containing it.  The certified outcome is:
\begin{enumerate}
\item the nodal subsystem has rank $11$ of $17$ and \emph{no coloop},
so, the two kernels agreeing, every exceptional curve of a small
resolution does appear with nonzero coefficient in some relation among
them: Friedman's necessary condition is satisfied;
\item the mixed matrix has $19$ rows on $21$ rays, its kernel has
dimension $6$, and \emph{both} $\dP_7$ coordinates vanish identically
on it, the root coordinate among them.
\end{enumerate}
By Theorem~\ref{thm:necessity}(2) and Lemma~\ref{lem:D}, $X_{20}$
admits no one-parameter smoothing.  Here the obstruction is carried
entirely by the divisorial direction, and a reader applying
\cite{Friedman86,RT09} alone would find nothing wrong.  On
$X^{\circ}$, by contrast, the two nodal coloops already signal that
the nodal theory is insufficient.

We also ran the test on the anticanonical hypersurface $X_{19}$ of the
$19$-vertex reflexive polytope $\Delta_{19}$ printed in
\cite[Section~5.3]{PaperIII}, whose generic member has fourteen nodes
and one $\dP_7$-cone point, all isolated.  The certified outcome is
structurally opposite to $X^{\circ}$:
\begin{enumerate}
\item the node subsystem has rank $12$ with two coloops: as for
$X^{\circ}$, no purely nodal mechanism can smooth $X_{19}$;
\item on the full kernel (dimension $3$) the complementary
$\dP_7$-coordinate $\alpha$ is already forced to zero, so every
relation class lies on the $(-2)$-line without imposing the branch
restriction; and
\item no covered coordinate is forced: a generic kernel class is
nonzero in all fourteen node coordinates and in the $(-2)$-coordinate
simultaneously.
\end{enumerate}
The necessity theory of this paper is therefore silent on $X_{19}$,
which motivates Question~\ref{q:X19}.  As a control in the opposite
direction, the test applied to the generic one-nodal quintic
threefold is silent as well, as it must be: there the exceptional
curve class vanishes in $H_2$ of the small resolution (the threefold
is $\QQ$-factorial), the kernel is the whole line $H_2(L;\QQ)$, and
the node coordinate is realized, the smoothable side of the nodal
picture.  A smoothing, if it exists,
must recruit the $\dP_7$ direction; the natural construction tools
are ambient toric degenerations in the sense of
Ilten--Vollmert~\cite{IV12}, and we leave the question to future
work.

\section{Discussion}\label{sec:discussion}

\subsection{The two-parameter degree-6 channel}
Theorem~\ref{thm:necessity} covers $\dP_6$ points smoothed through
the one-parameter component.  For the two-parameter (plane)
component the Milnor fiber has $b_3=2$, and the argument of
Lemma~\ref{lem:D} reduces the analogous nonvanishing to the
$\QQ$-linear independence of two periods of $\Omega_M$ on
$W\smallsetminus\Sb$, a finer question than nonvanishing, which
the toric method above does not decide.  None of the results of this
paper require it.

\subsection{First-order theory}
The tangent-level analogue of the test of
Section~\ref{sec:kernelcomp} identifies
$\ker(B^{\mathsf T})$ with the image of the localization map on
first-order deformations, by the local-to-global comparison of
Friedman--Laza~\cite{FL22a,FL22b}; from that viewpoint the content of
this paper is that the first-order forced-zero statements persist at
every order of contact, which no first-order argument can deliver.
The necessity statements of \cite{Friedman86,RT09} are of first-order
type, with the upgrade to actual smoothings supplied in the nodal
case by unobstructedness; the present method needs neither, which is
what allows it to reach the non-lci cone points.

\subsection{Sufficiency}
The converse direction, constructing mixed smoothings when the test
is silent, is open even for $X_{19}$.  The purely nodal comparison is
instructive: in dimension three the converse of Friedman's criterion
holds for projective nodal Calabi--Yau threefolds, by the first-order
relation together with the unobstructedness theorems of Kawamata and
Tian, while in higher odd dimensions Rollenske and Thomas prove a
converse for nodal hypersurfaces in projective space and explain why
a general converse should not be expected \cite[Section~5]{RT09}.  We record one
structural observation from the toric side: for a face $F$ with
$\ell(F^{\circ})=1$ the polar vertices adjacent to $F$ give ambient
deformation degrees whose slice polytopes are Minkowski-rigid, so any
ambient construction must use degrees on the extension of the dual
edge, the exact combinatorial locus of Altmann's local smoothing
parameter.  This will be taken up elsewhere.

\appendix

\section{Computational certificates}\label{app:cert}

The exact lattice and matrix claims used in the proofs are reproduced by
two short, self-contained Python programs in exact rational arithmetic,
independently of any computer-algebra system.  A third self-contained
program combines exact integral checks with a floating-point diagnostic
of the period parametrization, whose underlying identity is proved
exactly in Lemma~\ref{lem:D1}.  The general test program reuses the
companion repository's exact toric source modules.  The inputs are
Tables~\ref{tab:rays} and~\ref{tab:kperp}, the ancillary subdivision
and restriction data documents, and the polytope of
\cite[Section~5.3]{PaperIII}:
\begin{itemize}
\item the lattice package (\texttt{milnor\_kernels.py}): the boundary structure of the exceptional
surfaces, the norms and pairings of the $K^{\perp}$ rows, the root
decompositions, the $(-2)$-enumerations with rigorous search bounds,
and the isometries to the standard del Pezzo markings;
\item the matrix package (\texttt{global\_kernel.py}): the reconstruction of all $36$ rows from
the printed data, the ranks, coloops, kernels, the four profiles and
the enlarged profile with their forced-zero lists, and the
criticality statement of Section~\ref{sec:kernelcomp}(3);
\item the period package (\texttt{dp\_periods.py}): the exact fan combinatorics behind
\eqref{eq:omegaM} (the $\pm1$ pairings on all rays, the divisor
computation, the slice self-intersection patterns), the corner-chart
transitions, and the constancy of the pullback in
Lemma~\ref{lem:D1}, verified to machine precision at random points;
\item the repository-integrated test package (\texttt{examples.py}): the general implementation of
Section~\ref{sec:test}, validated end-to-end by reproducing every
number of Section~\ref{sec:kernelcomp} for $X^{\circ}$ from scratch,
and the $\Delta_{19}$ and $\Delta_{20}$ computations, including the
absence of nodal coloops on $\Delta_{20}$;
\item the mirror-partner package (\texttt{mirror\_partner.py}): the
reconstruction of the twenty-four rows of $B_{\Delta}$ from
Tables~\ref{tab:partnerrays} and~\ref{tab:partnersurf}, including the
derivation of both exceptional surface lattices from the face data of
Table~\ref{tab:partnerrays}, the reduced-rigidity of the
$\mathbb F_1$ germ by the zero-sum test on its edge vectors, the two
vanishing columns of item~(3), the root decomposition $A_1\oplus A_2$ of
$K^{\perp}(\dP_6)$, the absence of $(-2)$-classes in
$K^{\perp}(\mathbb F_1)$ and the square $-8$ of its generator, and the
ranks, coloops and forced-zero lists of Section~\ref{sec:partner}.
\end{itemize}
The fixed subdivision of Section~\ref{sec:matrix} and its
smoothness and projectivity are certified by an exact integral
support function; the $\Delta$-regularity of the reference equation
of $X^{\circ}$, including nondegeneracy along every face, is
certified at an integral witness (\texttt{resolution\_fan.sage} and
\texttt{global\_section.sage}, with five further Sage certificates
and the face-data module \texttt{fixed\_example.py}).  The
subdivision of Section~\ref{sec:partner}, its $94$ unimodular maximal
cones, the two stars and the twenty diagonals are certified the same
way by \texttt{mirror\_partner\_fan.sage}.

The arXiv ancillary bundle includes this directory together with the
shared exact toric modules imported by the repository-integrated
programs.  The same files are public, together with the programs and
data of the companion papers, at
\begin{center}
\url{https://github.com/bwuebben/calabi-yau-smoothability}\,.
\end{center}

\begin{thebibliography}{Wue26a}
\bibitem[Alt95]{Altmann95} K.~Altmann, \emph{Minkowski sums and homogeneous
deformations of toric varieties}, T\^ohoku Math.\ J.\ (2) \textbf{47}
(1995), 151--184.  (The statement cited here as (6.5) is (2.9), Case~2
of the preprint version, arXiv:alg-geom/9403003; it is restated as
\cite[(6.3)]{Altmann97}.)
\bibitem[Alt97]{Altmann97} K.~Altmann, \emph{The versal deformation of an
isolated toric Gorenstein singularity}, Invent.\ Math.\ \textbf{128}
(1997), 443--479.  (For the proof of the component description see also
arXiv:alg-geom/9403004v2, (8.2)--(8.3).)
\bibitem[Art68]{Artin68} M.~Artin, \emph{On the solutions of analytic
equations}, Invent.\ Math.\ \textbf{5} (1968), 277--291.
\bibitem[Art69]{Artin69} M.~Artin, \emph{Algebraic approximation of
structures over complete local rings}, Publ.\ Math.\ IH\'ES
\textbf{36} (1969), 23--58.
\bibitem[Bat94]{Batyrev94} V.~V. Batyrev, \emph{Dual polyhedra and mirror
symmetry for Calabi--Yau hypersurfaces in toric varieties}, J.\ Algebraic
Geom.\ \textbf{3} (1994), 493--535.
\bibitem[Con00]{Conrad00} B.~Conrad, \emph{Grothendieck duality and base
change}, Lecture Notes in Math.\ \textbf{1750}, Springer, 2000.
\bibitem[dJvS94]{dJvS94} T.~de Jong and D.~van Straten, \emph{On the
deformation theory of rational surface singularities with reduced
fundamental cycle}, J.\ Algebraic Geom.\ \textbf{3} (1994), 117--172.
\bibitem[Fri86]{Friedman86} R.~Friedman, \emph{Simultaneous resolution of
threefold double points}, Math.\ Ann.\ \textbf{274} (1986), 671--689.
\bibitem[FL24]{FL22b} R.~Friedman and R.~Laza, \emph{Deformations of some
local Calabi--Yau manifolds}, \'Epijournal G\'eom.\ Alg\'ebrique,
volume in honour of C.~Voisin (2024); arXiv:2203.11738.
\bibitem[FL25]{FL22a} R.~Friedman and R.~Laza, \emph{Deformations of
singular Fano and Calabi--Yau varieties}, J.\ Differential Geom.\
\textbf{131} (2025), 65--131.
\bibitem[Fuj80]{Fujita80a} T.~Fujita, \emph{On the structure of polarized
manifolds with total deficiency one, I}, J.\ Math.\ Soc.\ Japan
\textbf{32} (1980), 709--725.
\bibitem[Gra60]{Grauert60} H.~Grauert, \emph{Ein Theorem der analytischen
Garbentheorie und die Modulr\"aume komplexer Strukturen}, Publ.\ Math.\
IH\'ES \textbf{5} (1960), 5--64.
\bibitem[Gra72]{Grauert72} H.~Grauert, \emph{\"Uber die Deformation
isolierter Singularit\"aten analytischer Mengen}, Invent.\ Math.\
\textbf{15} (1972), 171--198.
\bibitem[Gra74]{Grauert74} H.~Grauert, \emph{Der Satz von Kuranishi f\"ur
kompakte komplexe R\"aume}, Invent.\ Math.\ \textbf{25} (1974), 107--142.
\bibitem[Gre20]{Greuel19} G.-M. Greuel, \emph{Deformation and smoothing of
singularities}, in: Handbook of Geometry and Topology of Singularities I,
Springer, 2020; arXiv:1903.03661.
\bibitem[Gro97]{Gross97} M.~Gross, \emph{Deforming Calabi--Yau threefolds},
Math.\ Ann.\ \textbf{308} (1997), 187--220.
\bibitem[IV12]{IV12} N.~Ilten and R.~Vollmert, \emph{Deformations of
rational $T$-varieties}, J.\ Algebraic Geom.\ \textbf{21} (2012),
473--493.
\bibitem[KP21]{KP21} A.-S. Kaloghiros and A.~Petracci, \emph{On toric
geometry and K-stability of Fano varieties}, Trans.\ Amer.\ Math.\
Soc.\ Ser.\ B \textbf{8} (2021), 548--577.
\bibitem[KS72]{KasSchlessinger72} A.~Kas and M.~Schlessinger, \emph{On the
versal deformation of a complex space with an isolated singularity},
Math.\ Ann.\ \textbf{196} (1972), 23--29.
\bibitem[KS00]{KS00} M.~Kreuzer and H.~Skarke, \emph{Complete
classification of reflexive polyhedra in four dimensions}, Adv.\ Theor.\
Math.\ Phys.\ \textbf{4} (2000), 1209--1230.
\bibitem[KP23]{KP23} A.~Kuznetsov and Yu.~Prokhorov, \emph{On
higher-dimensional del Pezzo varieties}, Izv.\ Math.\ \textbf{87} (2023),
no.~3, 488--561.
\bibitem[Mil68]{Milnor68} J.~Milnor, \emph{Singular points of complex
hypersurfaces}, Ann.\ of Math.\ Studies \textbf{61}, Princeton, 1968.
\bibitem[Pet19]{Petracci19} A.~Petracci, \emph{On deformations of toric
Fano varieties}, arXiv:1912.01538.
\bibitem[Pet20]{Petracci20} A.~Petracci, \emph{An example of mirror
symmetry for Fano threefolds}, in: Birational Geometry and Moduli Spaces,
Springer INdAM Ser.\ \textbf{39} (2020), 173--188.
\bibitem[Pin74]{Pinkham74} H.~Pinkham, \emph{Deformations of algebraic
varieties with $G_m$ action}, Ast\'erisque \textbf{20} (1974).
\bibitem[PP15]{PP15} P.~Popescu-Pampu, \emph{On the smoothings of
non-normal isolated surface singularities}, J.\ Singul.\ \textbf{12}
(2015), 164--179.
\bibitem[RT09]{RT09} S.~Rollenske and R.~Thomas, \emph{Smoothing nodal
Calabi--Yau $n$-folds}, J.\ Topol.\ \textbf{2} (2009), 405--421.
\bibitem[Wue26a]{PaperI} B.~J. Wuebben, \emph{Non-smoothable Calabi--Yau
threefolds from reflexive polytopes}, preprint, 2026; available at
\url{https://github.com/bwuebben/calabi-yau-smoothability}.
\bibitem[Wue26b]{PaperII} B.~J. Wuebben, \emph{Smoothing Calabi--Yau
threefolds in Gorenstein toric Fano fourfolds}, preprint, 2026; available
at \url{https://github.com/bwuebben/calabi-yau-smoothability}.
\bibitem[Wue26c]{PaperIII} B.~J. Wuebben, \emph{Doubly isolated Batyrev
mirror pairs and non-smoothable Calabi--Yau threefolds}, preprint, 2026;
available at \url{https://github.com/bwuebben/calabi-yau-smoothability}.
\end{thebibliography}

\bigskip
\noindent{\sc Bernd Johannes Wuebben, New York, NY,}
\texttt{wuebben@gmail.com}

\end{document}
